% Options for packages loaded elsewhere
% Options for packages loaded elsewhere
\PassOptionsToPackage{unicode}{hyperref}
\PassOptionsToPackage{hyphens}{url}
\PassOptionsToPackage{dvipsnames,svgnames,x11names}{xcolor}
%
\documentclass[
  11pt,
  letterpaper,
]{article}
\usepackage{xcolor}
\usepackage[margin=1in]{geometry}
\usepackage{amsmath,amssymb}
\setcounter{secnumdepth}{5}
\usepackage{iftex}
\ifPDFTeX
  \usepackage[T1]{fontenc}
  % \usepackage[utf8]{inputenc}
  \usepackage{textcomp} % provide euro and other symbols
\else % if luatex or xetex
  \usepackage{unicode-math} % this also loads fontspec
  \defaultfontfeatures{Scale=MatchLowercase}
  \defaultfontfeatures[\rmfamily]{Ligatures=TeX,Scale=1}
\fi
\usepackage{lmodern}
\ifPDFTeX\else
  % xetex/luatex font selection
\fi
% Use upquote if available, for straight quotes in verbatim environments
\IfFileExists{upquote.sty}{\usepackage{upquote}}{}
\IfFileExists{microtype.sty}{% use microtype if available
  \usepackage[]{microtype}
  \UseMicrotypeSet[protrusion]{basicmath} % disable protrusion for tt fonts
}{}
\makeatletter
\@ifundefined{KOMAClassName}{% if non-KOMA class
  \IfFileExists{parskip.sty}{%
    \usepackage{parskip}
  }{% else
    \setlength{\parindent}{0pt}
    \setlength{\parskip}{6pt plus 2pt minus 1pt}}
}{% if KOMA class
  \KOMAoptions{parskip=half}}
\makeatother
% Make \paragraph and \subparagraph free-standing
\makeatletter
\ifx\paragraph\undefined\else
  \let\oldparagraph\paragraph
  \renewcommand{\paragraph}{
    \@ifstar
      \xxxParagraphStar
      \xxxParagraphNoStar
  }
  \newcommand{\xxxParagraphStar}[1]{\oldparagraph*{#1}\mbox{}}
  \newcommand{\xxxParagraphNoStar}[1]{\oldparagraph{#1}\mbox{}}
\fi
\ifx\subparagraph\undefined\else
  \let\oldsubparagraph\subparagraph
  \renewcommand{\subparagraph}{
    \@ifstar
      \xxxSubParagraphStar
      \xxxSubParagraphNoStar
  }
  \newcommand{\xxxSubParagraphStar}[1]{\oldsubparagraph*{#1}\mbox{}}
  \newcommand{\xxxSubParagraphNoStar}[1]{\oldsubparagraph{#1}\mbox{}}
\fi
\makeatother


\usepackage{longtable,booktabs,array}
\newcounter{none} % for unnumbered tables
\usepackage{multirow}
\usepackage{calc} % for calculating minipage widths
% Correct order of tables after \paragraph or \subparagraph
\usepackage{etoolbox}
\makeatletter
\patchcmd\longtable{\par}{\if@noskipsec\mbox{}\fi\par}{}{}
\makeatother
% Allow footnotes in longtable head/foot
\IfFileExists{footnotehyper.sty}{\usepackage{footnotehyper}}{\usepackage{footnote}}
\makesavenoteenv{longtable}
\usepackage{graphicx}
\makeatletter
\newsavebox\pandoc@box
\newcommand*\pandocbounded[1]{% scales image to fit in text height/width
  \sbox\pandoc@box{#1}%
  \Gscale@div\@tempa{\textheight}{\dimexpr\ht\pandoc@box+\dp\pandoc@box\relax}%
  \Gscale@div\@tempb{\linewidth}{\wd\pandoc@box}%
  \ifdim\@tempb\p@<\@tempa\p@\let\@tempa\@tempb\fi% select the smaller of both
  \ifdim\@tempa\p@<\p@\scalebox{\@tempa}{\usebox\pandoc@box}%
  \else\usebox{\pandoc@box}%
  \fi%
}
% Set default figure placement to htbp
\def\fps@figure{htbp}
\makeatother


% definitions for citeproc citations
\NewDocumentCommand\citeproctext{}{}
\NewDocumentCommand\citeproc{mm}{%
  \begingroup\def\citeproctext{#2}\cite{#1}\endgroup}
\makeatletter
 % allow citations to break across lines
 \let\@cite@ofmt\@firstofone
 % avoid brackets around text for \cite:
 \def\@biblabel#1{}
 \def\@cite#1#2{{#1\if@tempswa , #2\fi}}
\makeatother
\newlength{\cslhangindent}
\setlength{\cslhangindent}{1.5em}
\newlength{\csllabelwidth}
\setlength{\csllabelwidth}{3em}
\newenvironment{CSLReferences}[2] % #1 hanging-indent, #2 entry-spacing
 {\begin{list}{}{%
  \setlength{\itemindent}{0pt}
  \setlength{\leftmargin}{0pt}
  \setlength{\parsep}{0pt}
  % turn on hanging indent if param 1 is 1
  \ifodd #1
   \setlength{\leftmargin}{\cslhangindent}
   \setlength{\itemindent}{-1\cslhangindent}
  \fi
  % set entry spacing
  \setlength{\itemsep}{#2\baselineskip}}}
 {\end{list}}
\usepackage{calc}
\newcommand{\CSLBlock}[1]{\hfill\break\parbox[t]{\linewidth}{\strut\ignorespaces#1\strut}}
\newcommand{\CSLLeftMargin}[1]{\parbox[t]{\csllabelwidth}{\strut#1\strut}}
\newcommand{\CSLRightInline}[1]{\parbox[t]{\linewidth - \csllabelwidth}{\strut#1\strut}}
\newcommand{\CSLIndent}[1]{\hspace{\cslhangindent}#1}



\setlength{\emergencystretch}{3em} % prevent overfull lines

\providecommand{\tightlist}{%
  \setlength{\itemsep}{0pt}\setlength{\parskip}{0pt}}



 


\DeclareUnicodeCharacter{2212}{\ensuremath{-}}
\DeclareUnicodeCharacter{2013}{--}
\DeclareUnicodeCharacter{2014}{---}
\DeclareUnicodeCharacter{2009}{\,}
\DeclareUnicodeCharacter{00A0}{~}
\DeclareUnicodeCharacter{00D7}{\ensuremath{\times}}
\DeclareUnicodeCharacter{2264}{\ensuremath{\leq}}
\DeclareUnicodeCharacter{2265}{\ensuremath{\geq}}
\DeclareUnicodeCharacter{2260}{\ensuremath{\neq}}
\DeclareUnicodeCharacter{2248}{\ensuremath{\approx}}
\DeclareUnicodeCharacter{03B1}{\ensuremath{\alpha}}
\DeclareUnicodeCharacter{03B2}{\ensuremath{\beta}}
\DeclareUnicodeCharacter{03C3}{\ensuremath{\sigma}}
\DeclareUnicodeCharacter{03C7}{\ensuremath{\chi}}
\DeclareUnicodeCharacter{2032}{\ensuremath{'}}
\DeclareUnicodeCharacter{2713}{\checkmark}

\usepackage{caption}
\captionsetup{labelsep=period}

\usepackage{multirow}
% \setcitestyle{round}

% \AtBeginEnvironment{table}{\footnotesize}
\AtBeginEnvironment{apptbl}{\footnotesize}
% Table font size: \footnotesize by default; table-size.lua emits \nexttablesize{\scriptsize}
% before a table whose div has size="scriptsize", which applies to that table only
\makeatletter
\newcommand\table@defaultsize{\footnotesize}
\let\table@size\table@defaultsize
\newcommand\nexttablesize[1]{\gdef\table@size{#1}}
\AtBeginEnvironment{longtable}{\table@size\global\let\table@size\table@defaultsize}
\makeatother

% Figure notes: note-in-float.lua wraps a figure's Note in \floatnote{...}, which
% prints it after the caption (the float package always sets the caption last)
\makeatletter
\let\floatnote@text\@empty
\newcommand\floatnote[1]{\gdef\floatnote@text{#1}}
\newcommand\fs@plainnote{\fs@plain
  \def\@fs@post{\ifx\floatnote@text\@empty\else
    \vspace\abovecaptionskip\noindent\floatnote@text\par
    \global\let\floatnote@text\@empty\fi}}
\makeatother

% Pin appendix floats in place so they stay next to their Note paragraphs
% (deferred inside \AtBeginDocument because Quarto's \newfloat comes after this file)
\AtBeginDocument{%
  \floatstyle{plainnote}\restylefloat{appfig}%
  \floatplacement{appfig}{H}\floatplacement{apptbl}{H}}

% Caption label for appendix tables typeset as breakable longtables (breakable-apptbl.lua);
% position=top restores the skip between the caption and the table's top rule
\AtBeginDocument{\providecommand*\apptblname{Table W}}
\captionsetup[apptbl]{position=top}

% ACL-style title block: bold title; one column per author (bold name, regular
% affiliation, small monospace email)
\makeatletter
\renewcommand\@maketitle{%
  \newpage\null\vskip 0.5em
  \begin{center}
    {\Large\bfseries \@title\par}
    \vskip 1.5em
    {\normalsize \@author\par}
  \end{center}
  \par\vskip 1em}
\newcommand\aclauthor[3]{% name, affiliation lines, email
  \begin{tabular}[t]{@{}c@{}}
    \textbf{#1}\\[1pt]
    #2\\[1pt]
    {\small\texttt{#3}}
  \end{tabular}}
\makeatother

% Hardcoded PDF authors. Runs at \begin{document}, so it overrides the plain
% \author{} Quarto generates from the YAML (the YAML authors still set PDF metadata).
\AtBeginDocument{%
  \author{%
    \aclauthor{Davood Wadi}
      {Desautels Faculty of Management\\ McGill University\\ Montreal, Quebec, Canada}
      {davood.wadi@hec.ca}%
    \hspace{3em}%
    \aclauthor{Yu Ma}
      {Desautels Faculty of Management\\ McGill University\\ Montreal, Quebec, Canada}
      {yu.ma@mcgill.ca}%
  }%
}
\makeatletter
\@ifpackageloaded{float}{}{\usepackage{float}}
\floatstyle{plain}
\@ifundefined{c@chapter}{\newfloat{apptbl}{h}{loapptbl}}{\newfloat{apptbl}{h}{loapptbl}[chapter]}
\floatname{apptbl}{Table W}
\floatstyle{plaintop}
\restylefloat{apptbl}
\newcommand*\quartoapptblref[1]{Table \hyperref[#1]{W\ref{#1}}}
\@ifpackageloaded{caption}{}{\usepackage{caption}}
\DeclareCaptionLabelFormat{quartoapptblreflabelformat}{#1#2}
\captionsetup[apptbl]{labelformat=quartoapptblreflabelformat}
\newcommand*\listofapptbls{\listof{apptbl}{List of Table Ws}}
\makeatother
\makeatletter
\@ifpackageloaded{float}{}{\usepackage{float}}
\floatstyle{plain}
\@ifundefined{c@chapter}{\newfloat{appfig}{h}{loappfig}}{\newfloat{appfig}{h}{loappfig}[chapter]}
\floatname{appfig}{Figure W}
\newcommand*\quartoappfigref[1]{Figure \hyperref[#1]{W\ref{#1}}}
\@ifpackageloaded{caption}{}{\usepackage{caption}}
\DeclareCaptionLabelFormat{quartoappfigreflabelformat}{#1#2}
\captionsetup[appfig]{labelformat=quartoappfigreflabelformat}
\newcommand*\listofappfigs{\listof{appfig}{List of Figure Ws}}
\makeatother
\makeatletter
\@ifpackageloaded{float}{}{\usepackage{float}}
\floatstyle{plain}
\@ifundefined{c@chapter}{\newfloat{appenv}{h}{loapp}}{\newfloat{appenv}{h}{loapp}[chapter]}
\floatname{appenv}{Web Appendix}
\floatstyle{plaintop}
\restylefloat{appenv}
\newcommand*\quartoappref[1]{Web \hyperref[#1]{Appendix\ref{#1}}}
\@ifpackageloaded{caption}{}{\usepackage{caption}}
\DeclareCaptionLabelFormat{quartoappreflabelformat}{#1#2}
\captionsetup[appenv]{labelformat=quartoappreflabelformat}
\newcommand*\listofappenvs{\listof{appenv}{List of Web Appendixs}}
\makeatother
\makeatletter
\@ifpackageloaded{caption}{}{\usepackage{caption}}
\AtBeginDocument{%
\ifdefined\contentsname
  \renewcommand*\contentsname{Table of contents}
\else
  \newcommand\contentsname{Table of contents}
\fi
\ifdefined\listfigurename
  \renewcommand*\listfigurename{List of Figures}
\else
  \newcommand\listfigurename{List of Figures}
\fi
\ifdefined\listtablename
  \renewcommand*\listtablename{List of Tables}
\else
  \newcommand\listtablename{List of Tables}
\fi
\ifdefined\figurename
  \renewcommand*\figurename{Figure}
\else
  \newcommand\figurename{Figure}
\fi
\ifdefined\tablename
  \renewcommand*\tablename{Table}
\else
  \newcommand\tablename{Table}
\fi
}
\@ifpackageloaded{float}{}{\usepackage{float}}
\floatstyle{ruled}
\@ifundefined{c@chapter}{\newfloat{codelisting}{h}{lop}}{\newfloat{codelisting}{h}{lop}[chapter]}
\floatname{codelisting}{Listing}
\newcommand*\listoflistings{\listof{codelisting}{List of Listings}}
\makeatother
\makeatletter
\makeatother
\makeatletter
\@ifpackageloaded{caption}{}{\usepackage{caption}}
\@ifpackageloaded{subcaption}{}{\usepackage{subcaption}}
\makeatother
\usepackage{bookmark}
\IfFileExists{xurl.sty}{\usepackage{xurl}}{} % add URL line breaks if available
\urlstyle{same}
% fallback for those not using the hyperref driver hyperxmp:
\makeatletter
\@ifundefined{xmpquote}{\newcommand{\xmpquote}[1]{#1}}{}
\makeatother
\hypersetup{
  pdftitle={Training decision models from choice microdata and published regression tables for silicon sampling},
  pdfauthor={Davood Wadi; Yu Ma},
  colorlinks=true,
  linkcolor={blue},
  filecolor={Maroon},
  citecolor={Blue},
  urlcolor={Blue},
  pdfcreator={LaTeX via pandoc}}


\title{Training decision models from choice microdata and published
regression tables for silicon sampling}
\author{Davood Wadi \and Yu Ma}
\date{}
\begin{document}
\maketitle


\textbf{Abstract}\\
Using off-the-shelf generative large language models (LLMs) to simulate
human respondents (silicon sampling) severely biases downstream
econometric inference, while standard fine-tuning requires individual
choice microdata that most published studies never release. We align
single-pass decision models that output normalized choice probabilities
directly over valid options without text generation, making linear
regression estimators directly differentiable in model weights. To
unlock studies that report only experimental designs and aggregate
regression tables, we derive an unbiased mini-batch
population-\(\boldsymbol{\beta}\) loss that pairs a frozen
base-probability control variate with split-sample cross-products,
stabilized by a Kullback-Leibler anchor on table-only data or
respondent-level cross-entropy (CE) on available microdata. We evaluate
open-weight decision models across five data-access regimes trained on
consumer delivery and electricity-plan choice studies and tested
zero-shot on an airline conjoint experiment and a six-arm
nutrition-labeling randomized trial. Microdata CE fine-tuning
substantially lowers out-of-sample treatment-effect mean squared error
(MSE), table-only training beats the off-the-shelf baseline without a
single respondent choice, and hybrid microdata-plus-table training
transfers missing price and green-energy trade-offs to close {[}TBD{]}\%
to {[}TBD{]}\% of the gap toward a full-microdata oracle.

\textbf{Keywords.} Decision models, Silicon sampling, Econometric
inference, Conjoint analysis, Marketing research

\section{Introduction}\label{sec-intro}

Simulating consumer choices, attitudes, and survey responses with
foundation models, commonly referred to as silicon sampling, has
attracted growing interest in marketing and behavioral research. In
marketing research, the primary goal is downstream econometric
inference, such as estimating attribute slopes in conjoint analysis or
causal treatment effects in randomized experiments. Currently, research
relies on off-the-shelf generative large language models (LLMs) that
generate text autoregressively. However, replacing human respondents
with off-the-shelf generative LLM outputs severely biases downstream
regressions because post-training for deterministic reasoning and safety
compliance collapses choice uncertainty and triggers safety refusals on
sensitive options (Goli and Singh 2024, Leng et al. 2026, Wang et al.
2026). Whereas existing post-inference methods adjust or reweight fixed
model outputs using target-study human data (Leng et al. 2026, Wang et
al. 2026), we align models before inference through their output
architecture and training supervision.

Architecturally, we adopt single-pass decision models (Palmer 2026).
Instead of generating text tokens, a decision model takes a respondent
persona and a structured choice task and outputs a normalized
probability distribution over the valid options in one forward pass.
This readout eliminates sampling variance and safety refusals while
making linear regression estimators, such as ordinary least squares
(OLS), directly differentiable in the model weights. Yet architecture
alone is insufficient because off-the-shelf decision models are
pretrained on single-label objective benchmarks, treating disagreement
among respondents who share a persona as classification error and
biasing regression estimates. Minimizing cross-entropy (CE) on
individual choice microdata resolves this miscalibration when respondent
choices are observed. In practice, however, raw choice microdata is
scarce, proprietary, or privacy-restricted, whereas thousands of
published marketing studies release only experimental designs and
aggregate regression tables.

Training a decision model directly on published regression coefficients
\(\boldsymbol{\beta}^*\) via mini-batch gradient descent presents two
statistical challenges. First, population OLS slopes depend on choice
probabilities across all tasks in a study rather than one mini-batch;
fitting OLS within small batches chases finite-sample noise, whereas
squaring a single-batch estimate of the population slopes introduces
variance bias that penalizes persona heterogeneity. Second, matching a
small vector of aggregate slopes leaves individual choice entropy
unconstrained, allowing predicted probabilities to drift toward zero or
one. We address the first challenge with an unbiased
population-\(\boldsymbol{\beta}\) loss that uses frozen base-model
probabilities as a control variate and multiplies estimation errors
across two disjoint half-batches without using respondent choices. We
address the second challenge by pairing this loss with an entropy
anchor, using Kullback-Leibler (KL) divergence toward base probabilities
for table-only training or respondent-level CE on available microdata
for hybrid training.

We evaluate this framework across a five-regime data-access ladder using
open-weight decision models (Palmer 2026) trained on a US drone-delivery
discrete-choice experiment (Berke et al. 2023) and an electricity-plan
conjoint study (Wu 2025), and evaluated zero-shot on an unseen
sustainable-aviation-fuel flight conjoint (Goldfarb et al. 2026) and a
six-arm nutrition-labeling randomized trial (Grummon et al. 2025). Our
design establishes three contributions. First, CE fine-tuning on
respondent microdata substantially lowers zero-shot treatment-effect
mean squared error (MSE) on unseen experiments relative to the
off-the-shelf model, whether training on single-study delivery microdata
or on full microdata across both studies as an upper-bound oracle.
Second, without access to respondent choices, training solely on the
eight published electricity slopes of Wu (2025) improves out-of-sample
MSE over the off-the-shelf baseline, and combining those published
slopes with Berke et al. (2023) delivery microdata transfers household
price and green-energy trade-offs that delivery microdata alone lacks.
Third, progressing from table-only to hybrid training closes {[}TBD{]}\%
to {[}TBD{]}\% of the out-of-sample MSE gap toward the unattainable
full-microdata oracle.

\section{Literature review}\label{sec-lit}

Early work in silicon sampling prompted generative LLMs with demographic
profiles to simulate human responses in political, economic, and
marketing surveys (Argyle et al. 2023, Brand et al. 2023, Horton et al.
2023, Santurkar et al. 2023, Durmus et al. 2024, Goli and Singh 2024,
Sarstedt et al. 2024). Although generative LLMs often capture the
direction of average preferences, prior studies document systematic
discrepancies from human populations, including mode collapse onto a
single option, attenuated within-group variance, demographic
stereotyping, and prompt sensitivity (Santurkar et al. 2023, Koo et al.
2024, Shu et al. 2024, Wadi and Fredette 2025, Deas and McKeown 2025,
Kabir et al. 2025, Zahradnickova et al. 2026, Jung et al. 2026). In
marketing applications, these distortions severely bias downstream
econometric estimation (Goli and Singh 2024, Leng et al. 2026, Wang et
al. 2026). In choice-based conjoint analysis, relying on simulated
choices or naively pooling them with human data biases preference
estimates, while scaling model capability or chain-of-thought reasoning
worsens error by inducing overly deterministic choices (Wang et al.
2026). Recent marketing studies address these discrepancies after
inference by treating the language model as a fixed generator, using
target-study human benchmark data to statistically debias downstream
regression parameters (Wang et al. 2026, Angelopoulos et al. 2023; AAE)
or optimally reweight ensembles of candidate prompts, personas, and
model checkpoints (Leng et al. 2026; HuMCal). Complementing these
post-inference methods, we align predictions inside the model before
inference through both output architecture and training supervision.

Architecturally, several studies avoid the cost and sampling variance of
repeated Monte Carlo text generation by extracting next-token
probabilities over option symbols (e.g., A, B, and C) from causal LLMs
({Kadavath et al.} 2022, Kuribayashi et al. 2024, Bentegeac et al. 2025,
Wadi et al. 2026). Yet symbol-level readouts remain sensitive to
vocabulary and token-frequency priors. Decision models, also known as
System One architectures, resolve this sensitivity via a single-pass
scoring head (Asseri and Asseri 2026, Deußer et al. 2026, Geuter et al.
2026, Li 2026, Li and Wagle 2026, Palmer 2026). This head evaluates each
supplied option directly in representation space and restricts
probability mass strictly to the valid options without open-ended
generation or safety refusals. Crucially for econometric training,
outputting exact choice probabilities in a single forward pass makes
linear regression operators over experimental designs directly
differentiable with respect to model parameters. However, existing
decision models are pretrained on single-label objective benchmarks such
as mathematics, code verification, and factual question answering,
causing severe probability miscalibration on subjective choice tasks
(Gao et al. 2026).

On the training dimension, existing fine-tuning approaches for silicon
sampling require access to unreleased response data. At the token level,
Wang et al. (2026) fine-tuned GPT-4o on 50 target-study conjoint tasks
under a generic persona to emit a single choice letter, finding that
small-sample single-token supervision worsened standalone parameter
error relative to the human benchmark. Recent studies improve
calibration by supervising causal language models or decision heads on
large-scale survey microdata and subgroup response histograms (Suh et
al. 2025, Ku et al. 2026, Huang et al. 2026). When respondent-level
choice microdata is available, minimizing CE on a probabilistic decision
head effectively calibrates conditional choice probabilities across
personas. Yet across conjoint analysis, discrete-choice modeling
(McFadden 1972, Train 2009), and randomized online or field trials,
published articles and archival industry reports rarely disclose raw
respondent choices, reporting instead experimental designs and estimated
regression tables. We bridge this data gap by supervising decision
models directly on published population OLS coefficients via an unbiased
split-sample moment loss, deployed either standalone with a
base-distribution KL anchor when no microdata exists or jointly
alongside respondent-level CE when augmenting in-house microdata with
external literature tables.

\section{Methodology}\label{sec-methodology}

\textbf{Decision model architecture and choice probabilities.} Consider
an experimental or survey study consisting of \(N\) discrete-choice
tasks indexed by \(t \in \{1, \dots, N\}\). Each task \(t\) pairs a
respondent's multi-attribute persona state \(\mathbf{s}_t\) (e.g., sex,
age, education, income, political party, and race) with a question stem
presenting \(K_t\) discrete answer options \(\{1, \dots, K_t\}\). When
individual choice microdata is available, let
\(\mathbf{y}_t = (y_{t,1}, \dots, y_{t,K_t})^\top \in \{0, 1\}^{K_t}\)
denote the observed one-hot choice vector on task \(t\), where
\(y_{t,k} = 1\) if the respondent selected option \(k\) and \(0\)
otherwise.

Given persona state \(\mathbf{s}_t\), a single-pass decision model
(Palmer 2026) scores all \(K_t\) options in one causal forward pass
without a vocabulary generation head. The prompt places \(\mathbf{s}_t\)
in a state segment, followed by the question stem, each option \(k\)
enclosed in dedicated option delimiters, and a terminal decision token
appended after the final option. Because the terminal decision token
follows all \(K_t\) options, its final-layer hidden state
\(\mathbf{h}^{\mathrm{dec}}_t(\mathbf{s}_t) \in \mathbb{R}^{d_{\mathrm{model}}}\)
causally attends to the persona, the question stem, and the full option
set. Let
\(\mathbf{h}_{t,k}(\mathbf{s}_t) \in \mathbb{R}^{d_{\mathrm{model}}}\)
denote the final-layer hidden state at the closing delimiter of option
\(k\). Using learned affine projections \((\mathbf{W}_q, \mathbf{b}_q)\)
and \((\mathbf{W}_k, \mathbf{b}_k)\) with
\(\mathbf{W}_q, \mathbf{W}_k \in \mathbb{R}^{d_{\mathrm{head}} \times d_{\mathrm{model}}}\)
and \(\mathbf{b}_q, \mathbf{b}_k \in \mathbb{R}^{d_{\mathrm{head}}}\),
the decision head computes a scalar score
\(z_{t,k}(\mathbf{s}_t; \theta)\) for each option \(k\):
\begin{equation}\protect\phantomsection\label{eq-pointer-head}{
z_{t,k}(\mathbf{s}_t; \theta) = \frac{\big(\mathbf{W}_q \mathbf{h}^{\mathrm{dec}}_t(\mathbf{s}_t) + \mathbf{b}_q\big)^\top \big(\mathbf{W}_k \mathbf{h}_{t,k}(\mathbf{s}_t) + \mathbf{b}_k\big)}{\sqrt{d_{\mathrm{head}}}}.
}\end{equation} Applying a softmax function yields the normalized choice
probability vector
\(\mathbf{p}_t(\theta) = (p_{t,1}(\theta), \dots, p_{t,K_t}(\theta))^\top\)
in a single pass:
\begin{equation}\protect\phantomsection\label{eq-model-pmf}{
p_{t,k}(\theta) = P_{\mathrm{model}}(Y_t = k \mid \mathbf{s}_t; \theta) = \frac{\exp\!\big(z_{t,k}(\mathbf{s}_t; \theta)\big)}{\sum_{m=1}^{K_t} \exp\!\big(z_{t,m}(\mathbf{s}_t; \theta)\big)}.
}\end{equation} When respondent choices \(\mathbf{y}_t\) are observed,
model parameters \(\theta\) can be supervised directly via mini-batch
cross-entropy over a sampled task batch \(S \subset \{1, \dots, N\}\):
\begin{equation}\protect\phantomsection\label{eq-resp-ce}{
\mathcal{L}_{\mathrm{CE}}(\theta; S) = -\frac{1}{|S|} \sum_{t \in S} \sum_{k=1}^{K_t} y_{t,k} \log p_{t,k}(\theta).
}\end{equation}

\textbf{Population OLS operators and published targets.} Concatenating
option probabilities across all \(N\) tasks in a study yields the flat
vector
\(\mathbf{p}(\theta) = (\mathbf{p}_1(\theta)^\top, \dots, \mathbf{p}_N(\theta)^\top)^\top \in [0, 1]^M\),
where \(M = \sum_{t=1}^N K_t\), with corresponding stacked human choices
\(\mathbf{y} \in \{0, 1\}^M\). Any linear econometric specification
estimated on the study can be written over \(R\) regression rows via a
full-column-rank design matrix
\(\mathbf{X} \in \mathbb{R}^{R \times c}\) (comprising an intercept, any
nuisance controls, and \(d\) focal treatment columns indexed by
\(\mathcal{T} \subset \{1, \dots, c\}\)) and a fixed sparse weight
matrix \(\mathbf{W} \in \mathbb{R}^{R \times M}\) that maps option
probabilities to regression outcomes
\(\mathbf{W}\mathbf{p}(\theta) \in \mathbb{R}^R\) (e.g., selecting each
displayed profile's choice probability in conjoint average marginal
component effects or the healthier item's choice probability in a
randomized trial). For training, dividing by the human outcome standard
deviation \(\sigma > 0\) places coefficients across studies on a common
dimensionless scale. Using the Moore-Penrose pseudoinverse
\(\mathbf{X}^\dagger = (\mathbf{X}^\top \mathbf{X})^{-1}\mathbf{X}^\top\),
define the fixed standardized population slope operator:
\begin{equation}\protect\phantomsection\label{eq-pop-operator}{
\mathbf{B} = \frac{1}{\sigma} \big(\mathbf{X}^\dagger\big)_{\mathcal{T}, :} \mathbf{W} \in \mathbb{R}^{d \times M},
}\end{equation} partitioned by task into column blocks
\(\mathbf{B} = [\mathbf{B}_1, \dots, \mathbf{B}_N]\) with
\(\mathbf{B}_t \in \mathbb{R}^{d \times K_t}\). Under
Equation~\ref{eq-pop-operator}, the model's standardized population OLS
slopes are an exact linear map of its single-pass probabilities,
\(\boldsymbol{\beta}(\theta) = \mathbf{B}\mathbf{p}(\theta) = \sum_{t=1}^N \mathbf{B}_t \mathbf{p}_t(\theta)\),
and a published regression table supplies the standardized human target
\(\boldsymbol{\beta}^* \in \mathbb{R}^d\) (which equals
\(\mathbf{B}\mathbf{y}\) without requiring individual choices
\(\mathbf{y}\)). When a study estimates \(G > 1\) regressions, we
construct \((\mathbf{B}^{(g)}, \boldsymbol{\beta}^{*(g)})\) for each
\(g \in \{1, \dots, G\}\) and average the resulting regression losses
equally across \(g\).

\textbf{Unbiased split-sample population-beta loss.} To align
model-implied slopes \(\boldsymbol{\beta}(\theta)\) with a published
table \(\boldsymbol{\beta}^*\), our target population objective is the
standardized slope MSE:
\begin{equation}\protect\phantomsection\label{eq-pop-mse}{
\mathcal{L}_{\mathrm{pop}}(\theta) = \frac{1}{d} \left\lVert \mathbf{B}\mathbf{p}(\theta) - \boldsymbol{\beta}^* \right\rVert_2^2.
}\end{equation} Evaluating Equation~\ref{eq-pop-mse} directly per
optimization step is intractable because
\(\mathbf{B}\mathbf{p}(\theta)\) sums over all \(N\) tasks whereas
stochastic gradient descent evaluates a mini-batch \(S\) of
\(|S| \ll N\) tasks. Naive mini-batch surrogates fail for two reasons.
Refitting OLS inside a small batch \(S\) requires unreleased choices
\(\mathbf{y}_t\) and squares ill-conditioned small-sample design
inverses, whereas substituting a single-batch estimate
\(\hat{\boldsymbol{\beta}}_S(\theta)\) into
\(\frac{1}{d}\lVert \hat{\boldsymbol{\beta}}_S(\theta) - \boldsymbol{\beta}^* \rVert_2^2\)
adds positive variance bias
\(\frac{1}{d}\mathrm{tr}\,\mathrm{Cov}(\hat{\boldsymbol{\beta}}_S(\theta))\),
which increases with across-persona heterogeneity in
\(\mathbf{p}_t(\theta)\) and penalizes the model for conditioning on
respondent personas.

To construct an unbiased, low-variance gradient estimator without
respondent choices or stale cached predictions, we run a single
no-gradient evaluation pass at initialization \(\theta_0\) to record
frozen base probabilities \(\mathbf{p}_0 = \mathbf{p}(\theta_0)\) and
base slopes \(\boldsymbol{\beta}_0 = \mathbf{B}\mathbf{p}_0\). Rewriting
\(\boldsymbol{\beta}(\theta) = \boldsymbol{\beta}_0 + \sum_{t=1}^N \mathbf{B}_t (\mathbf{p}_t(\theta) - \mathbf{p}_{0,t})\),
any uniform subsample \(S_m \subset \{1, \dots, N\}\) yields an exact
control-variate estimator of the current population slopes:
\begin{equation}\protect\phantomsection\label{eq-cv-beta}{
\hat{\boldsymbol{\beta}}_m(\theta) = \boldsymbol{\beta}_0 + \frac{N}{|S_m|} \sum_{t \in S_m} \mathbf{B}_t \big(\mathbf{p}_t(\theta) - \mathbf{p}_{0,t}\big).
}\end{equation} Because \(\mathbf{p}_0\) is fixed, \(\mathbf{p}_{0,t}\)
cancels in expectation at every step so
\(\mathbb{E}_{S_m}[\hat{\boldsymbol{\beta}}_m(\theta)] = \boldsymbol{\beta}(\theta)\)
without staleness bias, while subtracting \(\mathbf{p}_{0,t}\) reduces
mini-batch variance. Partitioning each shuffled batch \(S\)
(\(|S| \ge 2\)) into two disjoint halves \(S_1\) and \(S_2\) of sizes
\(\lfloor |S| / 2 \rfloor\) and \(\lceil |S| / 2 \rceil\), we define the
split-sample cross loss as the inner product of their estimation errors:
\begin{equation}\protect\phantomsection\label{eq-cross-loss}{
\mathcal{L}_{\beta\text{-cross}}(\theta; S_1, S_2) = \frac{1}{d} \big(\hat{\boldsymbol{\beta}}_1(\theta) - \boldsymbol{\beta}^*\big)^\top \big(\hat{\boldsymbol{\beta}}_2(\theta) - \boldsymbol{\beta}^*\big).
}\end{equation} Up to negligible \(\mathcal{O}(1/N)\) finite-population
sampling covariance without replacement,
\(\hat{\boldsymbol{\beta}}_1(\theta)\) and
\(\hat{\boldsymbol{\beta}}_2(\theta)\) are independent unbiased
estimators of \(\boldsymbol{\beta}(\theta)\), implying
\(\mathbb{E}[\mathcal{L}_{\beta\text{-cross}}(\theta; S_1, S_2)] = \mathcal{L}_{\mathrm{pop}}(\theta)\)
and
\(\mathbb{E}[\nabla_\theta \mathcal{L}_{\beta\text{-cross}}(\theta; S_1, S_2)] = \nabla_\theta \mathcal{L}_{\mathrm{pop}}(\theta)\)
(individual step realizations of Equation~\ref{eq-cross-loss} may be
negative).

\textbf{KL anchoring across training regimes.} Because \(d \ll M\),
minimizing Equation~\ref{eq-cross-loss} alone constrains only \(d\)
linear projections of \(\mathbf{p}(\theta)\) and leaves the
\((M - d)\)-dimensional nullspace \(\ker(\mathbf{B})\) unconstrained.
Without regularization along those unconstrained directions, prolonged
optimization can sharpen option logits and push task probabilities
toward zero or one while still maintaining
\(\mathbf{B}\mathbf{p}(\theta) \approx \boldsymbol{\beta}^*\),
distorting choice uncertainty and destabilizing out-of-sample
treatment-effect estimates. To stabilize table supervision without using
any respondent choices, we anchor predicted probabilities on table-study
tasks to the frozen base distribution \(\mathbf{p}_{0,t}\) through a
forward Kullback-Leibler (KL) divergence penalty:
\begin{equation}\protect\phantomsection\label{eq-kl-anchor}{
\mathcal{L}_{\mathrm{KL}}(\theta; S) = \frac{1}{|S|} \sum_{t \in S} \sum_{k=1}^{K_t} p_{0,t,k} \log \frac{p_{0,t,k}}{p_{t,k}(\theta)}.
}\end{equation} Combining Equation~\ref{eq-cross-loss} and
Equation~\ref{eq-kl-anchor} with slope weight \(\lambda > 0\) and anchor
weight \(\alpha \ge 0\) yields the table-study objective:
\begin{equation}\protect\phantomsection\label{eq-table-loss}{
\mathcal{L}_{\mathrm{table}}(\theta; S) = \lambda \, \mathcal{L}_{\beta\text{-cross}}(\theta; S_1, S_2) + \alpha \, \mathcal{L}_{\mathrm{KL}}(\theta; S).
}\end{equation} Evaluating Equation~\ref{eq-kl-anchor} requires zero
additional forward passes because \(\mathbf{p}_{0,t}\) is already stored
in memory for Equation~\ref{eq-cv-beta}. Furthermore,
\(\mathcal{L}_{\mathrm{KL}}(\theta; S)\) is convex in the option logits
with per-task logit gradient
\(\mathbf{p}_t(\theta) - \mathbf{p}_{0,t}\), which smoothly pulls
drifting probabilities back toward \(\mathbf{p}_{0,t}\) and prevents
collapse onto zero or one.

In the table-only regime, where no microdata is available for any study,
we minimize Equation~\ref{eq-table-loss} directly over mini-batches
drawn from the table study. In the hybrid regime, where in-house choice
microdata \(\mathcal{D}_{\mathrm{micro}}\) is paired with external
table-only studies \(\mathcal{D}_{\mathrm{table}}\), we interleave
single-study mini-batches across both sources and minimize
respondent-level CE on \(\mathcal{D}_{\mathrm{micro}}\) alongside the
same anchored table loss Equation~\ref{eq-table-loss} on
\(\mathcal{D}_{\mathrm{table}}\):
\begin{equation}\protect\phantomsection\label{eq-hybrid-loss}{
\mathcal{L}_{\mathrm{hybrid}}(\theta; S) = \mathbb{I}(S \subset \mathcal{D}_{\mathrm{micro}}) \, \mathcal{L}_{\mathrm{CE}}(\theta; S) + \mathbb{I}(S \subset \mathcal{D}_{\mathrm{table}}) \, \mathcal{L}_{\mathrm{table}}(\theta; S).
}\end{equation} Under Equation~\ref{eq-hybrid-loss},
\(\mathcal{L}_{\mathrm{CE}}\) calibrates choice uncertainty on the
microdata domain, \(\mathcal{L}_{\beta\text{-cross}}\) transfers missing
attribute trade-offs from external regression tables, and the KL anchor
directly regularizes the unconstrained nullspace of the external table
study's tasks while keeping table supervision uniform across regimes.
Setting \(\alpha = 0\) removes the KL anchor in either regime, enabling
a direct ablation of whether the anchor is needed when training on a
table alone or alongside microdata.

\section{Experimental setup}\label{sec-setup}

\textbf{Studies and split roles.} We evaluate the framework across five
published US choice studies spanning distinct domains with no domain
overlap between training and evaluation (Table~\ref{tbl-datasets}). Two
studies serve as training data. The first, \texttt{Electricity} (Wu
2025; 1,861 respondents, 18,591 choices), is a paired electricity-plan
conjoint experiment varying monthly bill (\texttt{\$105} to
\texttt{\$125}), renewable-energy share (\texttt{10\%} to
\texttt{100\%}), and local renewable production (\texttt{No} or
\texttt{Yes}), yielding \(d = 8\) OLS average marginal component effects
(AMCEs). Because Wu (2025) displays these eight AMCEs graphically in
Figure 3 rather than in a printed numerical table, we reproduce the
eight slopes using the study's published OLS specification and divide by
the outcome standard deviation \(\sigma\) to stand in for a published
regression table; in table-supervised regimes, all 18,591 respondent
choices are strictly withheld from the loss and gradients. The second
training study, \texttt{Drones} (Berke et al. 2023; 3,625 respondents,
36,297 choices), is a discrete-choice experiment on drone versus ground
delivery across four vendor contexts varying delivery fee (\texttt{\$0}
to \texttt{\$10}), wait time, and drone privacy, with no sustainability
attribute. In our design, \texttt{Drones} represents in-house respondent
microdata and enters training only through Equation~\ref{eq-resp-ce}.

Three unseen studies serve as out-of-domain benchmarks separated by both
split role and attribute overlap with training
(Table~\ref{tbl-datasets}). \texttt{Flights} (Goldfarb et al. 2026;
1,000 respondents, 5,000 choices, 12 OLS AMCEs) is a paired
airline-ticket conjoint experiment varying ticket price (\texttt{\$250}
to \texttt{\$325}), sustainable aviation fuel, airline, baggage,
legroom, and seat selection. We designate \texttt{Flights} as the
shared-attribute validation set, using it to select every training run's
checkpoint and to judge the KL-anchor ablation so that held-out test
sets never inform modeling choices. \texttt{Flights} therefore differs
from the held-out test sets both because it selects checkpoints and
because it shares the core price-versus-sustainability trade-off with
\texttt{Electricity} (and small delivery fees with \texttt{Drones}),
alongside four flight-specific attributes. The two strictly held-out
zero-shot benchmarks are disjoint-attribute tests whose attributes have
no overlap with training: \texttt{Food\ labels} (Grummon et al. 2025;
5,633 respondents, 33,798 choices, 5 OLS arm contrasts), a six-arm
front-of-package nutrition-labeling randomized trial measuring
healthier-product identification relative to the positive-label control,
and \texttt{Vaccines} (Kreps et al. 2020; 1,971 respondents, 9,855
choices, 11 OLS AMCEs), a three-option COVID-19 vaccine conjoint
experiment (two vaccine profiles plus an opt-out option) varying
efficacy, duration, side effects, FDA authorization, country of origin,
and endorsements. Neither \texttt{Food\ labels} nor \texttt{Vaccines} is
used during training, hyperparameter tuning, or checkpoint selection,
and neither includes price or sustainability attributes.

\begin{longtable}[]{@{}
  >{\raggedright\arraybackslash}p{(\linewidth - 14\tabcolsep) * \real{0.1250}}
  >{\raggedright\arraybackslash}p{(\linewidth - 14\tabcolsep) * \real{0.1250}}
  >{\raggedright\arraybackslash}p{(\linewidth - 14\tabcolsep) * \real{0.1250}}
  >{\raggedright\arraybackslash}p{(\linewidth - 14\tabcolsep) * \real{0.1250}}
  >{\raggedright\arraybackslash}p{(\linewidth - 14\tabcolsep) * \real{0.1250}}
  >{\raggedleft\arraybackslash}p{(\linewidth - 14\tabcolsep) * \real{0.1250}}
  >{\raggedleft\arraybackslash}p{(\linewidth - 14\tabcolsep) * \real{0.1250}}
  >{\raggedleft\arraybackslash}p{(\linewidth - 14\tabcolsep) * \real{0.1250}}@{}}
\caption{Summary of training, validation, and held-out test
datasets}\label{tbl-datasets}\tabularnewline
\toprule\noalign{}
\begin{minipage}[b]{\linewidth}\raggedright
Dataset
\end{minipage} & \begin{minipage}[b]{\linewidth}\raggedright
Study
\end{minipage} & \begin{minipage}[b]{\linewidth}\raggedright
Role
\end{minipage} & \begin{minipage}[b]{\linewidth}\raggedright
Shared attributes (price, sustainability)
\end{minipage} & \begin{minipage}[b]{\linewidth}\raggedright
Other study attributes
\end{minipage} & \begin{minipage}[b]{\linewidth}\raggedleft
Respondents
\end{minipage} & \begin{minipage}[b]{\linewidth}\raggedleft
Choices
\end{minipage} & \begin{minipage}[b]{\linewidth}\raggedleft
OLS slopes
\end{minipage} \\
\midrule\noalign{}
\endfirsthead
\toprule\noalign{}
\begin{minipage}[b]{\linewidth}\raggedright
Dataset
\end{minipage} & \begin{minipage}[b]{\linewidth}\raggedright
Study
\end{minipage} & \begin{minipage}[b]{\linewidth}\raggedright
Role
\end{minipage} & \begin{minipage}[b]{\linewidth}\raggedright
Shared attributes (price, sustainability)
\end{minipage} & \begin{minipage}[b]{\linewidth}\raggedright
Other study attributes
\end{minipage} & \begin{minipage}[b]{\linewidth}\raggedleft
Respondents
\end{minipage} & \begin{minipage}[b]{\linewidth}\raggedleft
Choices
\end{minipage} & \begin{minipage}[b]{\linewidth}\raggedleft
OLS slopes
\end{minipage} \\
\midrule\noalign{}
\endhead
\bottomrule\noalign{}
\endlastfoot
Electricity & Wu (2025) & Training (OLS slopes or microdata) & Monthly
bill (\texttt{\$105} to \texttt{\$125}), renewable share (\texttt{10\%}
to \texttt{100\%}), local renewables & None & 1,861 & 18,591 & 8 \\
Drones & Berke et al. (2023) & Training (microdata) & Delivery fee
(\texttt{\$0} to \texttt{\$10}); no sustainability & Wait time, drone
privacy & 3,625 & 36,297 & n/a \\
Flights† & Goldfarb et al. (2026) & Shared-attribute validation & Ticket
price (\texttt{\$250} to \texttt{\$325}), sustainable aviation fuel &
Airline, baggage, legroom, seat & 1,000 & 5,000 & 12 \\
Food labels & Grummon et al. (2025) & Disjoint-attribute test & None &
Front-of-package warning labels (6 arms) & 5,633 & 33,798 & 5 \\
Vaccines & Kreps et al. (2020) & Disjoint-attribute test & None &
Efficacy, duration, side effects, FDA, origin, endorsements & 1,971 &
9,855 & 11 \\
\end{longtable}

\emph{Note.} Shorthand labels in the first column are used throughout
all tables and figures. Shared attributes highlight price and
sustainability dimensions that overlap between training and evaluation,
whereas other study attributes are unique to each benchmark. Training
regimes append the supervision format in parentheses: microdata trains
on respondent choices via cross-entropy (Equation~\ref{eq-resp-ce}),
whereas OLS slopes trains on the eight published-specification OLS
slopes of Wu (2025) via Equation~\ref{eq-table-loss} while withholding
all 18,591 respondent choices. OLS slopes reports the number of
population regression slopes used in training or evaluation (n/a for
\texttt{Drones}, which enters training only as microdata). †
\texttt{Flights} (Goldfarb et al. 2026) is the shared-attribute
validation set used to select each run's checkpoint and to judge the
KL-anchor ablation.

\textbf{Data-access regimes.} Across every model scale, we compare six
main regimes spanning the data-access ladder. \texttt{Untrained}
evaluates the off-the-shelf decision model without fine-tuning (\(0\)
steps). \texttt{Electricity\ (OLS\ slopes)} trains solely on the eight
Wu (2025) OLS slopes using the anchored table objective
Equation~\ref{eq-table-loss} (\(\lambda = 1.0\), \(\alpha = 0.005\),
1,395 steps over \(3\) epochs) without using respondent choices, whereas
\texttt{Electricity\ (microdata)} trains via Equation~\ref{eq-resp-ce}
on all 18,591 Wu (2025) choices (1,395 steps) to provide a matched
within-study microdata reference. \texttt{Drones\ (microdata)} trains
via Equation~\ref{eq-resp-ce} on the 36,297 Berke et al. (2023) choices
(2,730 steps), representing a researcher's proprietary microdata
baseline. \texttt{Drones\ (microdata)\ +\ Electricity\ (OLS\ slopes)}
applies the anchored hybrid objective Equation~\ref{eq-hybrid-loss}
(\(\lambda = 1.0\), \(\alpha = 0.005\), 4,125 steps) to combine in-house
\texttt{Drones} microdata with external \texttt{Electricity} OLS slopes,
and \texttt{Drones\ +\ Electricity\ (microdata)} pools respondent
choices across both training studies under Equation~\ref{eq-resp-ce}
(4,125 steps) as a full-microdata reference. To isolate the effect of KL
anchoring, we additionally run unanchored counterparts (\(\alpha = 0\))
of both \texttt{Electricity\ (OLS\ slopes)} and
\texttt{Drones\ (microdata)\ +\ Electricity\ (OLS\ slopes)} under
identical optimization settings, together with a sensitivity sweep over
\(\alpha\) on \texttt{Electricity\ (OLS\ slopes)}, evaluated strictly on
the \texttt{Flights} validation set.

\textbf{Models, optimization, and evaluation metrics.} We evaluate the
open-weight \texttt{KEV} decision-model family (Gao et al. 2026, Palmer
2026) across three scales (\texttt{KEV-0.8B}, \texttt{KEV-4B}, and
\texttt{KEV-9B}). Each base checkpoint is loaded in \texttt{bfloat16}
precision and adapted with Low-Rank Adaptation (LoRA, rank \(r = 64\),
scaling \(\alpha_{\mathrm{LoRA}} = 128\), dropout \(0.05\)) across all
linear backbone projections alongside the trainable decision head. All
fine-tuned regimes use identical optimization settings across model
scales: AdamW with a peak learning rate of \(5 \times 10^{-6}\), a
OneCycle learning-rate schedule with a \(10\%\) warmup and cosine decay,
weight decay \(0.01\), gradient-norm clipping at \(1.0\), mini-batch
size \(|S| = 40\) tasks, random seed \(42\), and \(3\) epochs. Within
each epoch, tasks are partitioned into single-study mini-batches (and,
for \texttt{Drones}, single-vendor mini-batches) of \(40\) tasks and
interleaved in shuffled order, merging any one-task remainder on a table
study into another batch so every split cross-loss step satisfies
\(|S| \ge 2\). To compute exact gradients of
Equation~\ref{eq-table-loss} within GPU memory limits, each table-study
step uses a two-pass micro-chunked procedure (a first no-gradient pass
collects option logits across micro-chunks to evaluate
Equation~\ref{eq-table-loss} and its logit gradients, and a second pass
replays identical dropout states to backpropagate those gradients),
whereas microdata steps accumulate exact additive cross-entropy
gradients in a single pass.

Every run is validated on \texttt{Flights} every \(32\) optimization
steps (\(1,280\) training tasks), as well as at step \(0\) and the final
step, and the checkpoint attaining the lowest validation population
slope MSE is selected for test evaluation. Whereas training standardizes
table targets by \(\sigma\) in Equation~\ref{eq-pop-operator} and
Equation~\ref{eq-pop-mse}, all validation and held-out test metrics
compare unstandardized population OLS slopes
\(\tilde{\boldsymbol{\beta}}(\theta) = (\mathbf{X}^\dagger)_{\mathcal{T},:}\mathbf{W}\mathbf{p}(\theta)\)
against the corresponding human OLS slopes
\(\tilde{\boldsymbol{\beta}}^{\mathrm{human}} = (\mathbf{X}^\dagger)_{\mathcal{T},:}\mathbf{W}\mathbf{y}\)
in each study's native probability-point units. Our primary metric
across all benchmarks is unstandardized population slope MSE,
\(\frac{1}{d}\lVert \tilde{\boldsymbol{\beta}}(\theta) - \tilde{\boldsymbol{\beta}}^{\mathrm{human}} \rVert_2^2\)
(reported multiplied by \(10^3\)). As a secondary individual-level
diagnostic reported in the Web Appendix, we compute choice-distribution
KL divergence (Choice KL), defined as the mean negative natural log
probability
\(-\frac{1}{N}\sum_{t=1}^N \sum_{k=1}^{K_t} y_{t,k} \log p_{t,k}(\theta)\)
in nats across all tasks in a study.

\section{Study 1: Choice-Based Conjoint Analysis}\label{sec-study1}

To test whether fine-tuning on human response distributions generalizes
to out-of-sample human studies, we estimate each held-out study's
econometric model on the predicted responses and compare the resulting
regression parameters against the human estimates.

\subsection{Experimental Design and Multinomial Logit
Specification}\label{experimental-design-and-multinomial-logit-specification}

For our first held-out study, following the benchmark protocol of Wang
et al. (2026), we use a COVID-19 vaccine conjoint dataset (Kreps et al.
2020), summarized in Table~\ref{tbl-datasets}, to investigate part-worth
utility recovery. In this study, 1,971 US respondents each complete
\(t \in \{1, \dots, 5\}\) paired choice tasks (9,855 observations). Each
task presents two hypothetical vaccine profiles (\texttt{Option\ A} and
\texttt{Option\ B}) with binary attribute vectors
\(\mathbf{x}_{i,t,A}, \mathbf{x}_{i,t,B} \in \{0, 1\}^d\), alongside a
no-vaccine outside option (\texttt{Option\ C}) whose utility is
normalized to zero. Profiles vary across \(7\) attributes coded as
\(d = 11\) non-reference binary indicators
(\quartoapptblref{apptbl-wang-coefs}).

Under the standard Multinomial Logit (MNL) model (Wang et al. 2026),
choice probabilities for parameter vector
\(\boldsymbol{\beta} \in \mathbb{R}^d\) are:
\begin{equation}\protect\phantomsection\label{eq-mnl-probs}{
P_m(\mathbf{x}_{i,t}; \boldsymbol{\beta}) = \frac{\exp(\mathbf{x}_{i,t,m}^\top \boldsymbol{\beta})}{1 + \exp(\mathbf{x}_{i,t,A}^\top \boldsymbol{\beta}) + \exp(\mathbf{x}_{i,t,B}^\top \boldsymbol{\beta})}, \qquad m \in \{A, B\},
}\end{equation} with
\(P_C(\mathbf{x}_{i,t}; \boldsymbol{\beta}) = 1 - P_A(\mathbf{x}_{i,t}; \boldsymbol{\beta}) - P_B(\mathbf{x}_{i,t}; \boldsymbol{\beta})\).
Given observed or predicted choice shares
\((p_{i,t,A}, p_{i,t,B}, p_{i,t,C})\) across \(n\) choice sets, we
estimate \(\hat{\boldsymbol{\beta}}\) via Newton-CG maximum likelihood:
\begin{equation}\protect\phantomsection\label{eq-mnl-mle}{
\hat{\boldsymbol{\beta}} = \arg\min_{\boldsymbol{\beta} \in \mathbb{R}^d} -\frac{1}{n} \sum_{(i,t)} \sum_{m \in \{A, B, C\}} p_{i,t,m} \log P_m(\mathbf{x}_{i,t}; \boldsymbol{\beta}).
}\end{equation} Human respondents and generative LLMs supply one-hot
choices, whereas decision models supply continuous probabilities from
Equation~\ref{eq-model-pmf}. Let
\(\boldsymbol{\beta}^* \in \mathbb{R}^{11}\) denote the ground-truth MNL
vector estimated on all 9,855 human choices. Following Wang et al.
(2026), we evaluate models on a 200-respondent subsample (\emph{n} =
1,000 choice sets) using Mean Absolute Percentage Error
(\(\mathrm{MAPE}\)) and Mean Squared Error (\(\mathrm{MSE}\)):
\begin{equation}\protect\phantomsection\label{eq-wang-mape}{
\mathrm{MAPE}(\hat{\boldsymbol{\beta}}, \boldsymbol{\beta}^*) = \frac{1}{d} \sum_{j=1}^{d} \frac{|\hat{\beta}_j - \beta^*_j|}{|\beta^*_j|} \times 100\%,
}\end{equation}
\begin{equation}\protect\phantomsection\label{eq-wang-mse}{
\mathrm{MSE}(\hat{\boldsymbol{\beta}}, \boldsymbol{\beta}^*) = \frac{1}{d} \sum_{j=1}^{d} (\hat{\beta}_j - \beta^*_j)^2.
}\end{equation}

We test both unconditioned random personas and \(9\)-attribute
respondent personas (US region, gender, age, education, race, income,
religion, political party, and ideology;
\quartoapptblref{apptbl-attributes}). Decision models score the verbatim
survey instrument (Kreps et al. 2020), whereas generative LLMs
(\texttt{Gemini\ 3.5\ Flash-Lite} at \texttt{minimal} and \texttt{high}
thinking; \texttt{Gemini\ 4\ Argon} at \texttt{low} and \texttt{high}
thinking) use the Wang et al. (2026) prompt format (which additionally
states country). Unlike Wang et al. (2026), who excluded vaccine opt-out
respondents due to LLM safety refusals, we retain the full untruncated
respondent population and benchmark against MNL estimated directly on
the 200 sampled humans.

\subsection{Conjoint Estimation
Results}\label{conjoint-estimation-results}

\textbf{Generative LLM bias.} Estimating MNL on the 200 sampled humans
yields a finite-sample benchmark of \(\mathrm{MSE} = 0.0065\)
(Table~\ref{tbl-wang-mnl}; \quartoappfigref{appfig-wang-random} and
\quartoappfigref{appfig-wang-persona}). Consistent with Wang et al.
(2026), generative LLMs exhibit severe bias under random personas
(\(\mathrm{MSE} = 3.2073\) to \(26.2179\)). Conditioning on the \(9\)
respondent attributes cuts error substantially
(\(\mathrm{MSE} = 0.5925\) to \(2.1687\)), yet estimates remain far from
the human benchmark. Moreover, higher reasoning budgets and model scale
worsen parameter recovery. Under respondent personas, raising
\texttt{Gemini\ 3.5\ Flash-Lite} thinking from \texttt{minimal} to
\texttt{high} doubles MSE (\(0.5925\) to \(1.1984\)), while
\texttt{Gemini\ 4\ Argon} degrades further (\(1.7637\) at \texttt{low}
and \(2.1687\) at \texttt{high}).

\textbf{Why choice probabilities matter.} Generative LLMs bias conjoint
estimation because discrete text generation collapses choice uncertainty
and triggers safety refusals. Under random personas, generative LLMs
rarely select the no-vaccine option (\(0.0\%\) to \(3.0\%\)
vs.~\(20.3\%\) for humans), while \texttt{Gemini\ 4\ Argon\ (high)}
refuses \(16\) of \(1,000\) tasks. Even under respondent personas,
near-deterministic profile selection pushes conditional choice shares
toward zero or one, causing severe MNL scale inflation
(Figure~\ref{fig-wang-coefs}, Panel B;
\quartoapptblref{apptbl-wang-coefs}): \texttt{Efficacy:\ 90\%} reaches
\(+1.846\) to \(+3.538\) versus \(+0.746\) for humans. Scoring options
directly via \texttt{Kev\ 9B\ (base)} (\(T = 2.19\)) eliminates refusals
and cuts MSE by an order of magnitude (\(0.2449\) random; \(0.0815\)
respondent). Collapsing those same \texttt{Kev\ 9B\ (base)}
probabilities into one-hot top choices inflates respondent-persona MSE
nearly \(23\)-fold (\(0.0815\) to \(1.8678\)), showing that continuous
probability readouts--rather than discrete choices--are essential to
prevent scale distortion.

\textbf{Distributional survey fine-tuning.} Evaluated zero-shot--without
training or calibration on conjoint data--\texttt{Kev\ 9B\ (FT)} reduces
MNL MSE relative to \texttt{Kev\ 9B\ (base)} by \(64.3\%\) under random
personas (\(0.2449\) to \(0.0875\)) and by \(67.2\%\) under respondent
personas (\(0.0815\) to \(0.0267\); Table~\ref{tbl-wang-mnl}).
Fine-tuning also lowers choice cross-entropy (\texttt{1.6370} to
\texttt{1.5192} bits under random personas; \texttt{1.5255} to
\texttt{1.4449} bits under respondent personas) and raises predicted
no-vaccine shares (\(19.0\%\) and \(28.1\%\), up from \(11.9\%\) and
\(16.4\%\) for \texttt{Kev\ 9B\ (base)}). At the attribute level
(Figure~\ref{fig-wang-coefs}, Panel A;
\quartoapptblref{apptbl-wang-coefs}), \texttt{Kev\ 9B\ (base)}
overestimates the part-worths of vaccine efficacy and institutional
endorsements, with only \(5\) of \(11\) estimates falling within the
human subsample's \(95\%\) confidence intervals. By contrast,
\texttt{Kev\ 9B\ (FT)} recovers \(8\) of \(11\) estimates within those
\(95\%\) confidence intervals, closely aligning with full-sample human
part-worths on efficacy, FDA Emergency Use Authorization, country of
origin, and endorsements while attenuating protection duration and rare
major side effects toward zero.

\begin{longtable}[]{@{}
  >{\raggedright\arraybackslash}p{(\linewidth - 10\tabcolsep) * \real{0.2637}}
  >{\raggedright\arraybackslash}p{(\linewidth - 10\tabcolsep) * \real{0.2088}}
  >{\raggedleft\arraybackslash}p{(\linewidth - 10\tabcolsep) * \real{0.1209}}
  >{\raggedleft\arraybackslash}p{(\linewidth - 10\tabcolsep) * \real{0.1209}}
  >{\raggedleft\arraybackslash}p{(\linewidth - 10\tabcolsep) * \real{0.1209}}
  >{\raggedleft\arraybackslash}p{(\linewidth - 10\tabcolsep) * \real{0.1209}}@{}}
\caption{Multinomial logit part-worth parameter recovery on the COVID-19
vaccine conjoint benchmark}\label{tbl-wang-mnl}\tabularnewline
\toprule\noalign{}
\multirow{2}{=}{\begin{minipage}[b]{\linewidth}\raggedright
Estimator
\end{minipage}} &
\multirow{2}{=}{\begin{minipage}[b]{\linewidth}\raggedright
Configuration
\end{minipage}} &
\multicolumn{2}{>{\centering\arraybackslash}p{(\linewidth - 10\tabcolsep) * \real{0.2418} + 2\tabcolsep}}{%
\begin{minipage}[b]{\linewidth}\centering
Random Persona
\end{minipage}} &
\multicolumn{2}{>{\centering\arraybackslash}p{(\linewidth - 10\tabcolsep) * \real{0.2418} + 2\tabcolsep}@{}}{%
\begin{minipage}[b]{\linewidth}\centering
Respondent Persona
\end{minipage}} \\
& & \begin{minipage}[b]{\linewidth}\raggedleft
MAPE (\%)
\end{minipage} & \begin{minipage}[b]{\linewidth}\raggedleft
MSE
\end{minipage} & \begin{minipage}[b]{\linewidth}\raggedleft
MAPE (\%)
\end{minipage} & \begin{minipage}[b]{\linewidth}\raggedleft
MSE
\end{minipage} \\
\midrule\noalign{}
\endfirsthead
\toprule\noalign{}
\multirow{2}{=}{\begin{minipage}[b]{\linewidth}\raggedright
Estimator
\end{minipage}} &
\multirow{2}{=}{\begin{minipage}[b]{\linewidth}\raggedright
Configuration
\end{minipage}} &
\multicolumn{2}{>{\centering\arraybackslash}p{(\linewidth - 10\tabcolsep) * \real{0.2418} + 2\tabcolsep}}{%
\begin{minipage}[b]{\linewidth}\centering
Random Persona
\end{minipage}} &
\multicolumn{2}{>{\centering\arraybackslash}p{(\linewidth - 10\tabcolsep) * \real{0.2418} + 2\tabcolsep}@{}}{%
\begin{minipage}[b]{\linewidth}\centering
Respondent Persona
\end{minipage}} \\
& & \begin{minipage}[b]{\linewidth}\raggedleft
MAPE (\%)
\end{minipage} & \begin{minipage}[b]{\linewidth}\raggedleft
MSE
\end{minipage} & \begin{minipage}[b]{\linewidth}\raggedleft
MAPE (\%)
\end{minipage} & \begin{minipage}[b]{\linewidth}\raggedleft
MSE
\end{minipage} \\
\midrule\noalign{}
\endhead
\bottomrule\noalign{}
\endlastfoot
\multicolumn{6}{@{}>{\centering\arraybackslash}p{(\linewidth - 10\tabcolsep) * \real{0.9560} + 10\tabcolsep}@{}}{%
\emph{Human benchmark}} \\
Survey subsample & Observed choices & --- & --- & 22.1 & 0.0065 \\
\multicolumn{6}{@{}>{\centering\arraybackslash}p{(\linewidth - 10\tabcolsep) * \real{0.9560} + 10\tabcolsep}@{}}{%
\emph{Decision model}} \\
\multirow{2}{=}{Kev 9B} & Fine-tuned & 84.0 & 0.0875 & 56.4 & 0.0267 \\
& Off-the-shelf & 135.4 & 0.2449 & 93.5 & 0.0815 \\
\multicolumn{6}{@{}>{\centering\arraybackslash}p{(\linewidth - 10\tabcolsep) * \real{0.9560} + 10\tabcolsep}@{}}{%
\emph{Generative LLMs}} \\
\multirow{2}{=}{Gemini 3.5 Flash-Lite} & Minimal thinking & 531.7 &
3.2073 & 203.6 & 0.5925 \\
& High thinking & 998.2 & 14.4555 & 335.8 & 1.1984 \\
\multirow{2}{=}{Gemini 4 Argon} & Low thinking & 1,394.3 & 23.3998 &
376.7 & 1.7637 \\
& High thinking\textsuperscript{a} & 1,507.7 & 26.2179 & 395.1 &
2.1687 \\
\end{longtable}

\emph{Note.} Mean Absolute Percentage Error (MAPE) and Mean Squared
Error (MSE) are computed across \(d = 11\) vaccine attribute indicators
on the 200-respondent evaluation subsample (\emph{n} = 1,000 choice
sets), benchmarked against full-sample human \(\boldsymbol{\beta}^*\)
{[}\emph{N} = 1,971 respondents, 9,855 choice sets; Kreps et al. (2020);
Wang et al. (2026){]}. Under \texttt{Kev\ 9B}, \texttt{Fine-tuned}
(\texttt{FT}) uses unit temperature (\(T = 1.00\)) and
\texttt{Off-the-shelf} (\texttt{base}) uses released calibration
temperature (\(T = 2.19\)). \textsuperscript{a}Under random personas,
\emph{n} = 984 choice sets due to 16 safety refusals.

\begin{figure}

\centering{

\includegraphics[width=0.95\linewidth,height=\textheight,keepaspectratio]{figures/wang_mnl_coefs.png}

}

\caption{\label{fig-wang-coefs}Estimated multinomial logit part-worth
coefficients across COVID-19 vaccine attributes}

\vspace{\abovecaptionskip}

\noindent \emph{Note.} Estimates \(\hat{\boldsymbol{\beta}}\) and 95\%
asymptotic confidence intervals across the 11 non-reference vaccine
indicators under 9-attribute respondent personas (\emph{n} = 1,000
choice sets; generative LLM prompts additionally state country),
benchmarked against full-sample (\emph{N} = 9,855) and subsample
(\emph{n} = 1,000) human estimates (Kreps et al. 2020, Wang et al.
2026). Panel A compares \texttt{Kev\ 9B\ (base)} and
\texttt{Kev\ 9B\ (FT)} against human benchmarks; Panel B shows
coefficient inflation in generative LLMs across reasoning budgets.

\end{figure}%

\section{Study 2: Front-of-Package Nutrition Labeling
Trial}\label{sec-study2}

\subsection{Experimental Design and OLS
Specification}\label{experimental-design-and-ols-specification}

Our second held-out study evaluates causal treatment-effect recovery in
a front-of-package nutrition labeling trial (Grummon et al. 2025). In
this experiment, \emph{N} = 5,633 US grocery shoppers are randomized to
one of six between-subjects labeling arms (\texttt{Positive} serves as
the reference arm; \quartoapptblref{apptbl-grummon2025-coefs}). Each
shopper completes \(q \in \{1, \dots, 6\}\) three-option product choice
tasks across food and beverage categories, where \(y_1\) denotes the
objectively healthier option.

For each shopper \(i\), we compute the mean probability of selecting the
healthier product across the six tasks,
\(\bar{p}_i = \frac{1}{6} \sum_{q=1}^6 p_{i,q,1}\), where \(p_{i,q,1}\)
is the binary choice indicator \(y_{i,q,1}\) for human shoppers and the
predicted probability
\(P_{\mathrm{model}}(Y_{i,q} = y_1 \mid \mathbf{s}_i; \theta)\) under
the assigned arm for decision models. Following Grummon et al. (2025),
we regress \(\bar{p}_i\) on the five treatment-arm indicators via
ordinary least squares (OLS) with HC1 heteroskedasticity-robust standard
errors to obtain \(\hat{\boldsymbol{\beta}} \in \mathbb{R}^5\) relative
to \texttt{Positive}. We evaluate models on a 200-shopper subsample
(\emph{n} = 1,200 choices) against the full-sample human vector
\(\boldsymbol{\beta}^*\) (\emph{N} = 5,633). Each shopper profile
contains \(13\) attributes (\quartoapptblref{apptbl-attributes}): \(6\)
demographic attributes seen in survey training (sex, age, education,
income, race, and household size) and \(7\) covariates completely unseen
during training (Hispanic ethnicity, number of children, SNAP receipt,
WIC receipt, Nutrition Facts label usage, online grocery share, and
nutrition literacy score).

\subsection{Trial Estimation and Attribute Generalization
Results}\label{trial-estimation-and-attribute-generalization-results}

\textbf{Treatment-effect recovery.} Under \(13\)-attribute shopper
personas, scaling off-the-shelf models from \texttt{0.8B} to \texttt{9B}
worsens choice cross-entropy (\texttt{1.4770} to \texttt{1.5405} bits;
Table~\ref{tbl-grummon2025-ols}) and yields higher OLS MSE than
unconditioned random prompts at every scale
(\quartoappfigref{appfig-grummon2025-random};
\quartoappfigref{appfig-grummon2025-persona}). Survey fine-tuning
reverses this scaling pattern, steadily lowering cross-entropy from
\texttt{1.5536} bits (\texttt{0.8B}) to \texttt{1.3479} bits
(\texttt{9B}). Whereas smaller fine-tuned models (\texttt{0.8B} and
\texttt{4B}) improve arm rank ordering or MAPE without reducing OLS MSE,
\texttt{Kev\ 9B\ (FT)} improves every metric. It reaches a Spearman rank
correlation of \(\rho = +0.90\), lowers respondent-persona MAPE from
\(63.7\%\) to \(44.3\%\), and reduces OLS MSE by \(53.2\%\) (\(0.0047\)
to \(0.0022\); and by \(44.1\%\) to \(0.0019\) under random prompts),
approaching the 200-shopper human benchmark (\(\mathrm{MSE} = 0.0016\),
\(\mathrm{MAPE} = 22.7\%\)). At the arm level
(Figure~\ref{fig-grummon2025-coefs}, Panel A;
\quartoapptblref{apptbl-grummon2025-coefs}), \texttt{Kev\ 9B\ (FT)}
places all five estimates within the human subsample's \(95\%\)
confidence intervals, closely matching full-sample effects on the three
symmetrical interpretive labels (\texttt{Nutrition\ Info},
\texttt{Positive\ +\ Nutrition\ Info}, and \texttt{Spectrum}). Remaining
error concentrates in the two asymmetric \texttt{High\ In} arms, where
graphical warning shields in Grummon et al. (2025) were serialized into
text without visual warning cues.

\textbf{Generalization to unseen shopper attributes.} To test whether
fine-tuning merely memorizes training-corpus demographics, we evaluate
each model when prompted with one shopper attribute at a time across the
\(6\) seen demographics and the \(7\) unseen covariates
(Figure~\ref{fig-grummon2025-coefs}, Panel B;
\quartoapptblref{apptbl-grummon2025-personas}). Across individual
attributes, off-the-shelf \texttt{Kev\ 9B\ (base)} varies widely in
cross-entropy (\texttt{1.3935} to \texttt{1.6661} bits) and performs
worse than uniform guessing (\texttt{1.5850} bits) on
\texttt{Nutrition\ Facts\ label\ usage} (\texttt{1.6661}) and
\texttt{Share\ of\ groceries\ bought\ online} (\texttt{1.6016}). If
survey fine-tuning memorized only in-corpus demographics, improvements
would be confined to the \(6\) seen attributes. Instead,
\texttt{Kev\ 9B\ (FT)} lowers cross-entropy on every attribute,
converging to virtually identical mean cross-entropy on the \(7\) unseen
covariates (\texttt{1.3473} bits, vs.~\texttt{1.5328} for
\texttt{Kev\ 9B\ (base)}) and the \(6\) seen demographics
(\texttt{1.3477} bits, vs.~\texttt{1.4866} for
\texttt{Kev\ 9B\ (base)}).

\begin{longtable}[]{@{}
  >{\raggedright\arraybackslash}p{(\linewidth - 8\tabcolsep) * \real{0.2000}}
  >{\raggedleft\arraybackslash}p{(\linewidth - 8\tabcolsep) * \real{0.2000}}
  >{\raggedleft\arraybackslash}p{(\linewidth - 8\tabcolsep) * \real{0.2000}}
  >{\raggedleft\arraybackslash}p{(\linewidth - 8\tabcolsep) * \real{0.2000}}
  >{\raggedleft\arraybackslash}p{(\linewidth - 8\tabcolsep) * \real{0.2000}}@{}}
\caption{Choice cross-entropy and OLS treatment-effect recovery on the
nutrition labeling randomized
trial}\label{tbl-grummon2025-ols}\tabularnewline
\toprule\noalign{}
\begin{minipage}[b]{\linewidth}\raggedright
Estimator
\end{minipage} & \begin{minipage}[b]{\linewidth}\raggedleft
CE (bits)
\end{minipage} & \begin{minipage}[b]{\linewidth}\raggedleft
MAPE (\%)
\end{minipage} & \begin{minipage}[b]{\linewidth}\raggedleft
MSE
\end{minipage} & \begin{minipage}[b]{\linewidth}\raggedleft
Spearman \(\rho\)
\end{minipage} \\
\midrule\noalign{}
\endfirsthead
\toprule\noalign{}
\begin{minipage}[b]{\linewidth}\raggedright
Estimator
\end{minipage} & \begin{minipage}[b]{\linewidth}\raggedleft
CE (bits)
\end{minipage} & \begin{minipage}[b]{\linewidth}\raggedleft
MAPE (\%)
\end{minipage} & \begin{minipage}[b]{\linewidth}\raggedleft
MSE
\end{minipage} & \begin{minipage}[b]{\linewidth}\raggedleft
Spearman \(\rho\)
\end{minipage} \\
\midrule\noalign{}
\endhead
\bottomrule\noalign{}
\endlastfoot
Human subsample (\emph{N} = 200, \emph{n} = 1,200) & --- & 22.7\% &
0.0016 & +0.40 \\
Kev 0.8B (base) & 1.4770 & 58.5\% & 0.0074 & +0.50 \\
Kev 0.8B (FT) & 1.5536 & 66.5\% & 0.0078 & +0.70 \\
Kev 4B (base) & 1.5124 & 95.9\% & 0.0088 & +0.80 \\
Kev 4B (FT) & 1.4659 & 83.0\% & 0.0099 & +0.80 \\
Kev 9B (base) & 1.5405 & 63.7\% & 0.0047 & +0.80 \\
Kev 9B (FT) & 1.3479 & 44.3\% & 0.0022 & +0.90 \\
\end{longtable}

\emph{Note.} Three-option choice cross-entropy (CE, in bits; uniform
baseline \(\log_2 3 = 1.5850\) bits) and participant-level HC1 OLS
treatment-effect error across \(d = 5\) labeling arms (relative to
\texttt{Positive}) under 13-attribute respondent personas on the
evaluation subsample (\emph{N} = 200 shoppers, \emph{n} = 1,200
choices), benchmarked against full-sample human \(\boldsymbol{\beta}^*\)
(\emph{N} = 5,633, 33,798 choices; Grummon et al. (2025)). Off-the-shelf
(\texttt{base}) models use released calibration temperatures
(\(T = 2.35\), \(2.41\), and \(2.19\)); fine-tuned (\texttt{FT}) models
use unit temperature (\(T = 1.00\)).

\begin{figure}

\centering{

\includegraphics[width=0.95\linewidth,height=\textheight,keepaspectratio]{figures/grummon2025_ols_coefs.png}

}

\caption{\label{fig-grummon2025-coefs}Causal treatment-effect recovery
and shopper-attribute generalization on the nutrition labeling
randomized trial}

\vspace{\abovecaptionskip}

\noindent \emph{Note.} Evaluation on the 200-shopper subsample (\emph{n}
= 1,200 choices) of Grummon et al. (2025). Panel A compares
participant-level HC1 OLS treatment-effect estimates
\(\hat{\boldsymbol{\beta}}\) and 95\% confidence intervals across the 5
labeling arms (relative to \texttt{Positive}) under 13-attribute
respondent personas against full-sample (\emph{N} = 5,633) and subsample
(\emph{N} = 200) human benchmarks. Panel B compares single-attribute
choice cross-entropy (in bits, relative to uniform guessing at
\(\log_2 3 = 1.585\) bits) across 6 seen demographic attributes and 7
unseen shopper attributes.

\end{figure}%

\section{Discussion and Conclusion}\label{sec-discussion}

Our findings show that the severe econometric biases in generative LLM
silicon sampling (Goli and Singh 2024, Leng et al. 2026, Wang et al.
2026) stem from discrete text generation and single-label post-training
rather than fundamental limits of foundation models. Whereas recent
marketing methods correct fixed LLM outputs after inference using
target-study human data (Leng et al. 2026, Wang et al. 2026), aligning
models before inference requires pairing a probabilistic readout with
distributional supervision. Architecturally, single-pass decision heads
output normalized probabilities directly over valid options, eliminating
free-text parsing errors, Monte Carlo sampling variance, safety refusals
on sensitive opt-out choices, and the nearly \(23\)-fold MNL error
inflation caused by discretizing choices. Substantively, minimizing
cross-entropy against empirical sample distributions across 1.88 million
general survey responses (\texttt{US\ ISSP} and \texttt{Twin-2K-500})
calibrates conditional choice probabilities at unit temperature
(\(T = 1.00\)), generalizing zero-shot across unseen respondent
covariates and substantially improving conjoint MNL and experimental OLS
parameter recovery.

For marketing researchers, distributionally aligned decision models
enable fast, deterministic pre-testing before primary data collection
(e.g., screening candidate conjoint attributes, auditing survey
instruments, or evaluating policy messaging across demographic subgroups
in a single forward pass per task). At the same time, zero-shot
simulation complements rather than replaces human experimentation for
high-stakes managerial or regulatory decisions, as
\texttt{Kev\ 9B\ (FT)} still exhibits higher estimation error than a
200-participant human sample (\(\mathrm{MSE} = 0.0267\) vs.~\(0.0065\)
in conjoint MNL; \(0.0022\) vs.~\(0.0016\) in trial OLS). When
researchers do collect primary human responses in the target study,
upstream distributional alignment and downstream statistical calibration
(Leng et al. 2026, Wang et al. 2026) are naturally complementary.
Passing refusal-free, persona-conditioned probability vectors from
\texttt{Kev\ 9B\ (FT)} into AI-augmented estimators (Wang et al. 2026)
or checkpoint ensembles (Leng et al. 2026) could inject calibrated
respondent-level variation while reducing reliance on large human
calibration samples.

Three boundary conditions define priorities for future research. First,
our nutrition trial highlights a modality boundary of text-only silicon
sampling, as serializing graphical \texttt{High\ In} warning shields
into text loses visual salience cues. Extending distributional training
to multimodal vision-language decision heads that score package mockups
and advertisements directly is an immediate next step. Second, because
our training corpus centers on general US social and behavioral
attitudes, \texttt{Kev\ 9B\ (FT)} attenuates unfamiliar clinical
attributes (\texttt{Duration} and \texttt{Major\ side\ effects}) toward
zero, and smaller models (\texttt{0.8B} and \texttt{4B}) do not improve
OLS MSE. Incorporating historical commercial conjoint studies, scanner
panels, and cross-national surveys would broaden domain transfer. Third,
conditioning solely on static persona covariates \(\mathbf{s}_i\) treats
repeated choices per respondent as conditionally independent, motivating
sequential conditioning on prior choices or integration with mixed-logit
random-coefficients models. Together, single-pass decision architectures
and distributional survey training provide a practical, extensible
foundation for behavioral simulation in marketing science.

\section*{Declarations}\label{declarations}
\addcontentsline{toc}{section}{Declarations}

\subsection*{Funding and Competing
Interests}\label{funding-and-competing-interests}
\addcontentsline{toc}{subsection}{Funding and Competing Interests}

All authors certify that they have no affiliations with or involvement
in any organization or entity with any financial interest or
non-financial interest in the subject matter or materials discussed in
this manuscript. The authors have no funding to report.

\section*{References}\label{references}
\addcontentsline{toc}{section}{References}

\protect\phantomsection\label{refs}
\begin{CSLReferences}{1}{0}
\bibitem[\citeproctext]{ref-angelopoulos2023prediction}
Angelopoulos AN, Bates S, Fannjiang C, Jordan MI, Zrnic T (2023)
Prediction-powered inference. \emph{Science} 382(6671):669--674.

\bibitem[\citeproctext]{ref-argyle2023out}
Argyle LP, Busby EC, Fulda N, Gubler JR, Rytting C, Wingate D (2023) Out
of one, many: Using language models to simulate human samples.
\emph{Political Analysis} 31(3):337--351.

\bibitem[\citeproctext]{ref-asseri2026calibratedwhompersonalanguage}
Asseri B, Asseri A (2026)
\href{https://arxiv.org/abs/2609.36399}{Calibrated to whom? Persona and
language effects on cultural values in {JEV}}. \emph{arXiv preprint
arXiv:2609.36399}.

\bibitem[\citeproctext]{ref-bentegeac2025token}
Bentegeac R, Le Guellec B, Kuchcinski G, Amouyel P, Hamroun A (2025)
Token probabilities to mitigate large language models overconfidence in
answering medical questions: Quantitative study. \emph{Journal of
Medical Internet Research} 27:e64348.

\bibitem[\citeproctext]{ref-berke2023drone}
Berke A, Ding G, Chin C, Gopalakrishnan K, Larson K, Balakrishnan H, Li
MZ (2023) Drone delivery and the value of customer privacy: A discrete
choice experiment with US consumers. \emph{Transportation Research Part
C: Emerging Technologies} 157:104391.

\bibitem[\citeproctext]{ref-brand2023using}
Brand J, Israeli A, Ngwe D (2023) \emph{Using {LLMs} for market
research}. Working Paper 23-062 (Harvard Business School).

\bibitem[\citeproctext]{ref-deas-mckeown-2025-artificial}
Deas N, McKeown K (2025)
\href{https://doi.org/10.18653/v1/2025.emnlp-main.981}{Artificial
impressions: Evaluating large language model behavior through the lens
of trait impressions}. \emph{Proceedings of the 2025 conference on
{Empirical Methods in Natural Language Processing}}, 19418--19444.

\bibitem[\citeproctext]{ref-deusser2026evaluatingbenchmarkingmodeljev}
Deußer T, Sparrenberg L, Sifa R (2026)
\href{https://arxiv.org/abs/2609.37647}{Evaluating and benchmarking the
system one model {Jev}}. \emph{arXiv preprint arXiv:2609.37647}.

\bibitem[\citeproctext]{ref-durmus2024towards}
Durmus E, Nguyen K, Liao T, Schiefer N, Askell A, Bakhtin A, Chen C, et
al. (2024) \href{https://openreview.net/forum?id=zl16jLb91v}{Towards
measuring the representation of subjective global opinions in language
models}. \emph{First {Conference on Language Modeling}}.

\bibitem[\citeproctext]{ref-gao2026choicesfewerdecisionsordinalscale}
Gao T, Li J, Li Z, Chang Y, Wu Y (2026)
\href{https://arxiv.org/abs/2609.38827}{More choices, fewer decisions:
Ordinal-scale bias in {JEV}-like direct-decision models}. \emph{arXiv
preprint arXiv:2609.38827}.

\bibitem[\citeproctext]{ref-geuter2026releasinginnerjevsimple}
Geuter J, Cheng A, Alvarez-Melis D (2026)
\href{https://arxiv.org/abs/2610.17205}{Releasing the inner {Jev}: A
simple recipe for open decision models}. \emph{arXiv preprint
arXiv:2610.17205}.

\bibitem[\citeproctext]{ref-goldfarb2026comfort}
Goldfarb JL, Bhujle TR, Tester JW, Kriner DL (2026) Comfort or climate?
Consumer trade-offs between sustainable aviation fuel and flight
amenities. \emph{Journal of Cleaner Production} 577:149406.

\bibitem[\citeproctext]{ref-goli2024frontiers}
Goli A, Singh A (2024) Frontiers: Can large language models capture
human preferences? \emph{Marketing Science} 43(4):709--722.

\bibitem[\citeproctext]{ref-grummon2025nutrition}
Grummon AH, O'Sullivan K, Petimar J, Lee CJY, Zeitlin AB, Cleveland LP,
Collis C, et al. (2025) Nutrition info and other front-of-package labels
and simulated food and beverage purchases: A randomized clinical trial.
\emph{JAMA Network Open} 8(10):e2537389.
\url{https://doi.org/10.1001/jamanetworkopen.2025.37389}.

\bibitem[\citeproctext]{ref-horton2023large}
Horton JJ, Filippas A, Manning BS (2023) \emph{Large language models as
simulated economic agents: What can we learn from {Homo Silicus}?}
Working Paper 31122 (National Bureau of Economic Research).

\bibitem[\citeproctext]{ref-huang-etal-2026-distribution}
Huang J, Li M, Shao S (2026)
\href{https://doi.org/10.18653/v1/2026.findings-acl.457}{Distribution
shift alignment helps {LLM}s simulate survey response distributions}.
\emph{Findings of the {A}ssociation for {C}omputational {L}inguistics:
{ACL} 2026}, 9395--9409.

\bibitem[\citeproctext]{ref-jung-etal-2026-psychometric}
Jung J, Lutz M, Sen I, Strohmaier M (2026)
\href{https://doi.org/10.18653/v1/2026.eacl-long.380}{Do psychometric
tests work for large language models? Evaluation of tests on sexism,
racism, and morality}. \emph{Proceedings of the 19th conference of the
{European Chapter of the Association for Computational Linguistics}
(volume 1: Long papers)}, 8143--8173.

\bibitem[\citeproctext]{ref-kabir-etal-2025-break}
Kabir M, Abrar A, Ananiadou S (2025)
\href{https://doi.org/10.18653/v1/2025.emnlp-main.2}{Break the checkbox:
Challenging closed-style evaluations of cultural alignment in {LLM}s}.
\emph{Proceedings of the 2025 conference on {Empirical Methods in
Natural Language Processing}}, 24--51.

\bibitem[\citeproctext]{ref-kadavath2022language}
{Kadavath S, Conerly T, Askell A, Henighan T, Drain D, Perez E, Schiefer
N, et al.} (2022) Language models (mostly) know what they know.
\emph{arXiv preprint arXiv:2207.05221}.

\bibitem[\citeproctext]{ref-koo2024benchmarking}
Koo R, Lee M, Raheja V, Park JI, Kim ZM, Kang D (2024) Benchmarking
cognitive biases in large language models as evaluators. \emph{Findings
of the {Association for Computational Linguistics}: {ACL} 2024},
517--545.

\bibitem[\citeproctext]{ref-kreps2020factors}
Kreps S, Prasad S, Brownstein JS, Hswen Y, Garibaldi BT, Zhang B, Kriner
DL (2020) Factors associated with {US} adults' likelihood of accepting
COVID-19 vaccination. \emph{JAMA Network Open} 3(10):e2025594.

\bibitem[\citeproctext]{ref-ku2026siliconsamplingcrosssurveytransfer}
Ku CT, Hsu C, Huang PC, Liu FC, Cheng IL, Kang Y (2026)
\href{https://arxiv.org/abs/2607.03091}{Silicon sampling via
cross-survey transfer}. \emph{arXiv preprint arXiv:2607.03091}.

\bibitem[\citeproctext]{ref-kuribayashi-etal-2024-psychometric}
Kuribayashi T, Oseki Y, Baldwin T (2024)
\href{https://doi.org/10.18653/v1/2024.findings-naacl.129}{Psychometric
predictive power of large language models}. \emph{Findings of the
{Association for Computational Linguistics}: {NAACL} 2024}, 1983--2005.

\bibitem[\citeproctext]{ref-leng2026reducing}
Leng Y, Sang Y, Agarwal A (2026) Reducing disparity between {LLM}s and
humans: Optimal {LLM} sample calibration. \emph{Marketing Science}.
\url{https://doi.org/10.1287/mksc.2024.0983}.

\bibitem[\citeproctext]{ref-li2026kitescalingjevpopulation}
Li H (2026) \href{https://arxiv.org/abs/2609.27535}{{KITE}: Scaling
{Jev} population experiments with sparse flagship calibration}.
\emph{arXiv preprint arXiv:2609.27535}.

\bibitem[\citeproctext]{ref-li2026llmasjevllmsjevstyledecision}
Li Y, Wagle J (2026)
\href{https://arxiv.org/abs/2610.02076}{{LLM-as-Jev}: {LLMs} are already
{Jev}-style decision models -- when and how to fine-tune them}.
\emph{arXiv preprint arXiv:2610.02076}.

\bibitem[\citeproctext]{ref-mcfadden1972conditional}
McFadden D (1972) Conditional logit analysis of qualitative choice
behavior.

\bibitem[\citeproctext]{ref-palmer2026kev}
Palmer J (2026) Kev: Small {Jev}-like decision models.
\url{https://github.com/jaredpalmer/kev}.

\bibitem[\citeproctext]{ref-pmlr-v202-santurkar23a}
Santurkar S, Durmus E, Ladhak F, Lee C, Liang P, Hashimoto T (2023)
\href{https://proceedings.mlr.press/v202/santurkar23a.html}{Whose
opinions do language models reflect?} \emph{Proceedings of the 40th
{International Conference on Machine Learning}}, 29971--30004.

\bibitem[\citeproctext]{ref-sarstedt2024using}
Sarstedt M, Adler SJ, Rau L, Schmitt B (2024) Using large language
models to generate silicon samples in consumer and marketing research:
Challenges, opportunities, and guidelines. \emph{Psychology \&
Marketing} 41(6):1254--1270.

\bibitem[\citeproctext]{ref-shu-etal-2024-dont}
Shu B, Zhang L, Choi M, Dunagan L, Logeswaran L, Lee M, Card D, Jurgens
D (2024) \href{https://doi.org/10.18653/v1/2024.naacl-long.295}{You
don{'}t need a personality test to know these models are unreliable:
Assessing the reliability of large language models on psychometric
instruments}. \emph{Proceedings of the 2024 conference of the {North
American Chapter of the Association for Computational Linguistics}:
{Human Language Technologies} (volume 1: Long papers)}, 5263--5281.

\bibitem[\citeproctext]{ref-suh-etal-2025-language}
Suh J, Jahanparast E, Moon S, Kang M, Chang S (2025)
\href{https://doi.org/10.18653/v1/2025.acl-long.1028}{Language model
fine-tuning on scaled survey data for predicting distributions of public
opinions}. \emph{Proceedings of the 63rd annual meeting of the
{A}ssociation for {C}omputational {L}inguistics (volume 1: Long
papers)}, 21147--21170.

\bibitem[\citeproctext]{ref-train2009discrete}
Train KE (2009) \emph{Discrete choice methods with simulation}
(Cambridge university press).

\bibitem[\citeproctext]{ref-wadi2025monte}
Wadi D, Fredette M (2025) A {Monte-Carlo} sampling framework for
reliable evaluation of large language models using behavioral analysis.
\emph{Findings of the {Association for Computational Linguistics}:
{EMNLP} 2025}, 9414--9432.

\bibitem[\citeproctext]{ref-wadi2026every}
Wadi D, Ghodrat M, Philp M (2026) Every token counts: Exact
{Likert}-scale distributions for measuring LLM attitudes and biases.
\emph{arXiv preprint arXiv:2608.10503}.

\bibitem[\citeproctext]{ref-wang2026large}
Wang M, Zhang DJ, Zhang H (2026) Large language models for market
research: A data-augmentation approach. \emph{Marketing Science}
45(4):728--751. \url{https://doi.org/10.1287/mksc.2025.0009}.

\bibitem[\citeproctext]{ref-wu2025partisan}
Wu VY (2025) Partisan energy preferences in the united states:
Republicans prioritize price, democrats also consider renewables.
\emph{Energy Research \& Social Science} 123:103997.

\bibitem[\citeproctext]{ref-zahradnickova2026ai}
Zahradnickova K, Polisensky J, Lukeš M, Hranický R, Pouč P (2026) Which
AI? A comparative evaluation of generative AI and neuro-symbolic AI for
silicon sampling. \emph{Psychology \& Marketing}.
\url{https://doi.org/10.1002/mar.70272}.

\end{CSLReferences}

\newpage{}

\section*{Web Appendix}\label{web-appendix}
\addcontentsline{toc}{section}{Web Appendix}

{\def\LTcaptype{apptbl}
\begin{longtable}[]{@{}
  >{\raggedright\arraybackslash}p{(\linewidth - 14\tabcolsep) * \real{0.1250}}
  >{\raggedright\arraybackslash}p{(\linewidth - 14\tabcolsep) * \real{0.1250}}
  >{\raggedleft\arraybackslash}p{(\linewidth - 14\tabcolsep) * \real{0.1250}}
  >{\raggedleft\arraybackslash}p{(\linewidth - 14\tabcolsep) * \real{0.1250}}
  >{\raggedleft\arraybackslash}p{(\linewidth - 14\tabcolsep) * \real{0.1250}}
  >{\raggedleft\arraybackslash}p{(\linewidth - 14\tabcolsep) * \real{0.1250}}
  >{\raggedleft\arraybackslash}p{(\linewidth - 14\tabcolsep) * \real{0.1250}}
  >{\raggedleft\arraybackslash}p{(\linewidth - 14\tabcolsep) * \real{0.1250}}@{}}
\caption{{\label{apptbl-issp-waves}}Summary of the 10 US waves of the
International Social Survey Programme (2015 to 2024)}\tabularnewline
\toprule\noalign{}
\begin{minipage}[b]{\linewidth}\raggedright
Wave
\end{minipage} & \begin{minipage}[b]{\linewidth}\raggedright
Thematic module
\end{minipage} & \begin{minipage}[b]{\linewidth}\raggedleft
Questions
\end{minipage} & \begin{minipage}[b]{\linewidth}\raggedleft
Likert (\texttt{score})
\end{minipage} & \begin{minipage}[b]{\linewidth}\raggedleft
Categorical (\texttt{choice})
\end{minipage} & \begin{minipage}[b]{\linewidth}\raggedleft
Respondents
\end{minipage} & \begin{minipage}[b]{\linewidth}\raggedleft
Mean Qs / respondent
\end{minipage} & \begin{minipage}[b]{\linewidth}\raggedleft
Responses
\end{minipage} \\
\midrule\noalign{}
\endfirsthead
\toprule\noalign{}
\begin{minipage}[b]{\linewidth}\raggedright
Wave
\end{minipage} & \begin{minipage}[b]{\linewidth}\raggedright
Thematic module
\end{minipage} & \begin{minipage}[b]{\linewidth}\raggedleft
Questions
\end{minipage} & \begin{minipage}[b]{\linewidth}\raggedleft
Likert (\texttt{score})
\end{minipage} & \begin{minipage}[b]{\linewidth}\raggedleft
Categorical (\texttt{choice})
\end{minipage} & \begin{minipage}[b]{\linewidth}\raggedleft
Respondents
\end{minipage} & \begin{minipage}[b]{\linewidth}\raggedleft
Mean Qs / respondent
\end{minipage} & \begin{minipage}[b]{\linewidth}\raggedleft
Responses
\end{minipage} \\
\midrule\noalign{}
\endhead
\bottomrule\noalign{}
\endlastfoot
2015 & Work orientations IV & 82 & 68 & 14 & 1,474 & 49.1 & 72,416 \\
2016 & Role of government V & 60 & 52 & 8 & 1,384 & 58.6 & 81,043 \\
2017 & Social networks and social resources & 60 & 32 & 28 & 1,173 &
59.2 & 69,500 \\
2018 & Religion IV & 62 & 50 & 12 & 1,175 & 60.2 & 70,719 \\
2019 & Social inequality V & 49 & 44 & 5 & 1,852 & 45.9 & 84,963 \\
2020 & Environment IV & 55 & 45 & 10 & 1,847 & 52.3 & 96,598 \\
2021 & Health and health care II & 58 & 50 & 8 & 1,146 & 54.8 &
62,793 \\
2022 & Family and changing gender roles V & 55 & 32 & 23 & 1,126 & 45.6
& 51,303 \\
2023 & National identity and citizenship & 60 & 56 & 4 & 1,589 & 56.3 &
89,391 \\
2024 & Digital societies & 75 & 45 & 30 & 1,590 & 55.8 & 88,758 \\
Total & All 10 US waves (2015 to 2024) & 616 & 474 & 142 & 14,356 & 53.5
& 767,484 \\
\end{longtable}
}

\emph{Note.} Counts combine training and held-out validation items
extracted from the US subsample of each annual ISSP thematic module.
Four ordinal Likert (\texttt{score}) items from the 2020 Environment IV
wave (7,011 responses across 1,827 respondents) are held out for
validation; all remaining 612 items (760,473 responses) enter training.

{\def\LTcaptype{apptbl}
\begin{longtable}[]{@{}
  >{\raggedright\arraybackslash}p{(\linewidth - 4\tabcolsep) * \real{0.3333}}
  >{\raggedright\arraybackslash}p{(\linewidth - 4\tabcolsep) * \real{0.3333}}
  >{\raggedright\arraybackslash}p{(\linewidth - 4\tabcolsep) * \real{0.3333}}@{}}
\caption{{\label{apptbl-attributes}}Harmonized respondent persona
attributes across training and evaluation datasets}\tabularnewline
\toprule\noalign{}
\begin{minipage}[b]{\linewidth}\raggedright
Prompt attribute
\end{minipage} & \begin{minipage}[b]{\linewidth}\raggedright
Harmonized levels
\end{minipage} & \begin{minipage}[b]{\linewidth}\raggedright
Datasets
\end{minipage} \\
\midrule\noalign{}
\endfirsthead
\toprule\noalign{}
\begin{minipage}[b]{\linewidth}\raggedright
Prompt attribute
\end{minipage} & \begin{minipage}[b]{\linewidth}\raggedright
Harmonized levels
\end{minipage} & \begin{minipage}[b]{\linewidth}\raggedright
Datasets
\end{minipage} \\
\midrule\noalign{}
\endhead
\bottomrule\noalign{}
\endlastfoot
Sex & 3 (Male, Female, Non-binary) & ISSP, Twin-2K, Conjoint, Nutrition
RCT \\
Age & 4 (18-29, 30-49, 50-64, 65+) & ISSP, Twin-2K, Conjoint, Nutrition
RCT \\
Education & 4 (Less than high school, High school or some college,
Associate's degree, College graduate/some postgrad) & ISSP, Twin-2K,
Conjoint, Nutrition RCT \\
Religion & 12 (Protestant, Catholic, Orthodox, \ldots, Other) & ISSP,
Twin-2K, Conjoint \\
Attendance of religious services & 6 (More than once a week, Once a
week, Once or twice a month, A few times a year, Seldom, Never) & ISSP,
Twin-2K \\
Political ideology & 8 (Very liberal, Liberal, Somewhat liberal, \ldots,
Other) & ISSP, Twin-2K, Conjoint \\
Voted in the last general election & 2 (Yes, No) & ISSP \\
Place of living & 5 (A big city, The suburbs or outskirts of a big city,
A town or a small city, A country village, A farm or home in the
country) & ISSP \\
Marital status & 6 (Married, Living with a partner, Divorced, Separated,
Widowed, Never been married) & ISSP, Twin-2K \\
Main activity & 10 (Full-time employment, Part-time employment,
Self-employed, \ldots, Other) & ISSP, Twin-2K \\
Trade union membership & 2 (Yes, No) & ISSP \\
Political party & 5 (Democrat, Republican, Independent, Something else,
Other candidate) & ISSP, Twin-2K, Conjoint \\
Income & 15 (Less than \$30,000, \$30,000 - \$50,000, \$50,000 -
\$75,000, \ldots, \$35,000 - \$49,999) & ISSP, Twin-2K, Conjoint,
Nutrition RCT \\
Region & 4 (Northeast, Midwest, South, West) & ISSP, Twin-2K,
Conjoint \\
Ancestry & 96 (Algeria, American Indian, American only, \ldots,
Yugoslavia) & ISSP \\
Race & 5 (White, Black, Asian, Hispanic, Other) & Twin-2K, Conjoint,
Nutrition RCT \\
How many people, including yourself, live in your household? & 5 (1, 2,
3, 4, More than 4) & Twin-2K, Nutrition RCT \\
Hispanic or Latino & 2 (No, Yes) & Nutrition RCT \\
How many children (ages 0-18) currently live in your household? & 4 (0,
1, 2, 3 or more) & Nutrition RCT \\
SNAP / Food Stamps recipient & 3 (Don't know, No, Yes) & Nutrition
RCT \\
Do you or anyone else in your household currently receive WIC benefits?
& 3 (Don't know, No, Yes) & Nutrition RCT \\
Nutrition Facts label usage & 5 (Always, Most of the time, Never,
Rarely, Sometimes) & Nutrition RCT \\
Share of groceries bought online & 4 (25-49\%, 50-74\%, 75\% or more,
Less than 25\%) & Nutrition RCT \\
Nutrition literacy score & 5 (0 out of 4, 1 out of 4, 2 out of 4, 3 out
of 4, 4 out of 4) & Nutrition RCT \\
\end{longtable}
}

\emph{Note.} Attribute levels show the harmonized natural-language
categories rendered in each respondent's persona state \(\mathbf{s}_i\)
(abbreviated with ellipses where category lists exceed six values). The
bottom seven attributes appear exclusively in the held-out nutrition
labeling randomized clinical trial (Grummon et al. 2025) and are never
observed during survey training.

\nexttablesize{\scriptsize}
{\def\LTcaptype{apptbl}
\begin{longtable}[]{@{}
  >{\raggedright\arraybackslash}p{(\linewidth - 16\tabcolsep) * \real{0.1111}}
  >{\raggedleft\arraybackslash}p{(\linewidth - 16\tabcolsep) * \real{0.1111}}
  >{\raggedleft\arraybackslash}p{(\linewidth - 16\tabcolsep) * \real{0.1111}}
  >{\raggedleft\arraybackslash}p{(\linewidth - 16\tabcolsep) * \real{0.1111}}
  >{\raggedleft\arraybackslash}p{(\linewidth - 16\tabcolsep) * \real{0.1111}}
  >{\raggedleft\arraybackslash}p{(\linewidth - 16\tabcolsep) * \real{0.1111}}
  >{\raggedleft\arraybackslash}p{(\linewidth - 16\tabcolsep) * \real{0.1111}}
  >{\raggedleft\arraybackslash}p{(\linewidth - 16\tabcolsep) * \real{0.1111}}
  >{\raggedleft\arraybackslash}p{(\linewidth - 16\tabcolsep) * \real{0.1111}}@{}}
\caption{{\label{apptbl-wang-coefs}}Estimated multinomial logit
part-worth coefficients on the COVID-19 vaccine conjoint benchmark}\tabularnewline
\toprule\noalign{}
\begin{minipage}[b]{\linewidth}\raggedright
Vaccine attribute level
\end{minipage} & \begin{minipage}[b]{\linewidth}\raggedleft
Full human (\emph{N} = 9,855)
\end{minipage} & \begin{minipage}[b]{\linewidth}\raggedleft
Human subsample (\emph{n} = 1,000)
\end{minipage} & \begin{minipage}[b]{\linewidth}\raggedleft
Kev 9B (FT)
\end{minipage} & \begin{minipage}[b]{\linewidth}\raggedleft
Kev 9B (base)
\end{minipage} & \begin{minipage}[b]{\linewidth}\raggedleft
Gemini 3.5 Flash-Lite (high)
\end{minipage} & \begin{minipage}[b]{\linewidth}\raggedleft
Gemini 3.5 Flash-Lite (minimal)
\end{minipage} & \begin{minipage}[b]{\linewidth}\raggedleft
Gemini 4 Argon (high)
\end{minipage} & \begin{minipage}[b]{\linewidth}\raggedleft
Gemini 4 Argon (low)
\end{minipage} \\
\midrule\noalign{}
\endfirsthead
\toprule\noalign{}
\begin{minipage}[b]{\linewidth}\raggedright
Vaccine attribute level
\end{minipage} & \begin{minipage}[b]{\linewidth}\raggedleft
Full human (\emph{N} = 9,855)
\end{minipage} & \begin{minipage}[b]{\linewidth}\raggedleft
Human subsample (\emph{n} = 1,000)
\end{minipage} & \begin{minipage}[b]{\linewidth}\raggedleft
Kev 9B (FT)
\end{minipage} & \begin{minipage}[b]{\linewidth}\raggedleft
Kev 9B (base)
\end{minipage} & \begin{minipage}[b]{\linewidth}\raggedleft
Gemini 3.5 Flash-Lite (high)
\end{minipage} & \begin{minipage}[b]{\linewidth}\raggedleft
Gemini 3.5 Flash-Lite (minimal)
\end{minipage} & \begin{minipage}[b]{\linewidth}\raggedleft
Gemini 4 Argon (high)
\end{minipage} & \begin{minipage}[b]{\linewidth}\raggedleft
Gemini 4 Argon (low)
\end{minipage} \\
\midrule\noalign{}
\endhead
\bottomrule\noalign{}
\endlastfoot
Efficacy: 70\% (vs.~50\%) & +0.361 (0.035) & +0.280 (0.112) & +0.217
(0.113) & +0.565 (0.115) & +1.344 (0.158) & +0.688 (0.135) & +1.940
(0.176) & +1.671 (0.163) \\
Efficacy: 90\% (vs.~50\%) & +0.746 (0.035) & +0.679 (0.110) & +0.562
(0.111) & +1.074 (0.116) & +2.764 (0.188) & +1.846 (0.145) & +3.538
(0.218) & +3.148 (0.196) \\
Duration: 5 years (vs.~1 year) & +0.261 (0.029) & +0.277 (0.092) &
+0.035 (0.093) & +0.105 (0.095) & +0.940 (0.131) & +0.592 (0.114) &
+0.948 (0.140) & +0.830 (0.133) \\
Major side effects: 1 in 1,000,000 (vs.~1 in 10,000) & +0.304 (0.029) &
+0.389 (0.091) & +0.043 (0.092) & +0.152 (0.093) & +2.353 (0.154) &
+1.103 (0.114) & +1.647 (0.146) & +1.634 (0.140) \\
Minor side effects: 1 in 30 (vs.~1 in 10) & +0.076 (0.029) & +0.085
(0.090) & -0.043 (0.092) & +0.024 (0.093) & +0.341 (0.124) & +0.051
(0.109) & +0.282 (0.128) & +0.389 (0.124) \\
FDA: Emergency Use Authorization (vs.~Full approval) & -0.115 (0.029) &
-0.113 (0.092) & -0.123 (0.093) & -0.068 (0.095) & -1.089 (0.136) &
-0.978 (0.115) & -1.223 (0.141) & -1.329 (0.137) \\
Origin: United Kingdom (vs.~United States) & -0.140 (0.035) & -0.026
(0.109) & -0.192 (0.110) & -0.173 (0.113) & -0.234 (0.151) & -0.402
(0.127) & -0.527 (0.157) & -0.710 (0.151) \\
Origin: China (vs.~United States) & -0.570 (0.036) & -0.425 (0.113) &
-0.553 (0.114) & -0.624 (0.118) & -1.908 (0.179) & -2.449 (0.164) &
-3.556 (0.227) & -3.341 (0.210) \\
Endorsed: Joe Biden (vs.~Donald Trump) & +0.095 (0.040) & +0.100 (0.123)
& +0.135 (0.126) & +0.441 (0.128) & +0.600 (0.171) & +0.325 (0.148) &
+0.531 (0.178) & +0.385 (0.170) \\
Endorsed: US CDC (vs.~Donald Trump) & +0.416 (0.039) & +0.287 (0.123) &
+0.621 (0.124) & +1.017 (0.127) & +1.045 (0.172) & +0.635 (0.148) &
+1.175 (0.179) & +0.784 (0.169) \\
Endorsed: WHO (vs.~Donald Trump) & +0.273 (0.040) & +0.327 (0.124) &
+0.516 (0.126) & +0.730 (0.129) & +0.166 (0.173) & +0.334 (0.152) &
+0.699 (0.181) & +0.411 (0.173) \\
\end{longtable}
}

\emph{Note.} Point estimates and asymptotic standard errors (in
parentheses) across the 11 binary vaccine attribute indicators on the
200-respondent subsample (\emph{n} = 1,000 choice sets) under
respondent-persona conditioning, alongside full-sample human estimates
(\emph{N} = 1,971 respondents, 9,855 choice sets; Kreps et al. (2020);
Wang et al. (2026)). Respondent personas comprise 9 demographic and
political attributes (generative LLM prompts additionally state
country). Kev rows use the question wording, option labels, and
attribute descriptions of the original survey instrument (Kreps et al.
2020); the generative LLMs are prompted with the Wang et al. (2026)
wording.

\begin{appfig}

\centering{

\includegraphics[width=0.9\linewidth,height=\textheight,keepaspectratio]{figures/wang_mnl_random.png}

}

\caption{\label{appfig-wang-random}Multinomial logit part-worth
estimation error under unconditioned random-persona prompts}

\floatnote{

\emph{Note.} Panel A reports Mean Absolute Percentage Error (MAPE) and
Panel B reports Mean Squared Error (MSE) across the \(d = 11\) vaccine
attribute levels on the 200-respondent COVID-19 conjoint subsample
(\emph{n} = 1,000 choice sets; \emph{n} = 984 for
\texttt{Gemini\ 4\ Argon\ (high)} after 16 neutrality refusals) relative
to full-sample human \(\boldsymbol{\beta}^*\) (\emph{N} = 9,855; Kreps
et al. (2020); Wang et al. (2026)).

}

\end{appfig}%

\begin{appfig}

\centering{

\includegraphics[width=0.9\linewidth,height=\textheight,keepaspectratio]{figures/wang_mnl_persona.png}

}

\caption{\label{appfig-wang-persona}Multinomial logit part-worth
estimation error under respondent-persona conditioning}

\floatnote{

\emph{Note.} Panel A reports Mean Absolute Percentage Error (MAPE) and
Panel B reports Mean Squared Error (MSE) across the \(d = 11\) vaccine
attribute levels on the 200-respondent COVID-19 conjoint subsample
(\emph{n} = 1,000 choice sets) under 9-attribute respondent personas
(generative LLM prompts additionally state country) relative to
full-sample human \(\boldsymbol{\beta}^*\) (\emph{N} = 9,855; Kreps et
al. (2020); Wang et al. (2026)).

}

\end{appfig}%

\nexttablesize{\scriptsize}
{\def\LTcaptype{apptbl}
\begin{longtable}[]{@{}
  >{\raggedright\arraybackslash}p{(\linewidth - 16\tabcolsep) * \real{0.1111}}
  >{\raggedleft\arraybackslash}p{(\linewidth - 16\tabcolsep) * \real{0.1111}}
  >{\raggedleft\arraybackslash}p{(\linewidth - 16\tabcolsep) * \real{0.1111}}
  >{\raggedleft\arraybackslash}p{(\linewidth - 16\tabcolsep) * \real{0.1111}}
  >{\raggedleft\arraybackslash}p{(\linewidth - 16\tabcolsep) * \real{0.1111}}
  >{\raggedleft\arraybackslash}p{(\linewidth - 16\tabcolsep) * \real{0.1111}}
  >{\raggedleft\arraybackslash}p{(\linewidth - 16\tabcolsep) * \real{0.1111}}
  >{\raggedleft\arraybackslash}p{(\linewidth - 16\tabcolsep) * \real{0.1111}}
  >{\raggedleft\arraybackslash}p{(\linewidth - 16\tabcolsep) * \real{0.1111}}@{}}
\caption{{\label{apptbl-grummon2025-coefs}}Participant-level OLS
treatment-effect estimates on the nutrition labeling randomized trial}\tabularnewline
\toprule\noalign{}
\begin{minipage}[b]{\linewidth}\raggedright
Trial arm (vs.~Positive)
\end{minipage} & \begin{minipage}[b]{\linewidth}\raggedleft
Full human (\emph{N} = 5,633)
\end{minipage} & \begin{minipage}[b]{\linewidth}\raggedleft
Human subsample (\emph{N} = 200)
\end{minipage} & \begin{minipage}[b]{\linewidth}\raggedleft
Kev 0.8B (base)
\end{minipage} & \begin{minipage}[b]{\linewidth}\raggedleft
Kev 0.8B (FT)
\end{minipage} & \begin{minipage}[b]{\linewidth}\raggedleft
Kev 4B (base)
\end{minipage} & \begin{minipage}[b]{\linewidth}\raggedleft
Kev 4B (FT)
\end{minipage} & \begin{minipage}[b]{\linewidth}\raggedleft
Kev 9B (base)
\end{minipage} & \begin{minipage}[b]{\linewidth}\raggedleft
Kev 9B (FT)
\end{minipage} \\
\midrule\noalign{}
\endfirsthead
\toprule\noalign{}
\begin{minipage}[b]{\linewidth}\raggedright
Trial arm (vs.~Positive)
\end{minipage} & \begin{minipage}[b]{\linewidth}\raggedleft
Full human (\emph{N} = 5,633)
\end{minipage} & \begin{minipage}[b]{\linewidth}\raggedleft
Human subsample (\emph{N} = 200)
\end{minipage} & \begin{minipage}[b]{\linewidth}\raggedleft
Kev 0.8B (base)
\end{minipage} & \begin{minipage}[b]{\linewidth}\raggedleft
Kev 0.8B (FT)
\end{minipage} & \begin{minipage}[b]{\linewidth}\raggedleft
Kev 4B (base)
\end{minipage} & \begin{minipage}[b]{\linewidth}\raggedleft
Kev 4B (FT)
\end{minipage} & \begin{minipage}[b]{\linewidth}\raggedleft
Kev 9B (base)
\end{minipage} & \begin{minipage}[b]{\linewidth}\raggedleft
Kev 9B (FT)
\end{minipage} \\
\midrule\noalign{}
\endhead
\bottomrule\noalign{}
\endlastfoot
Nutrition Info label (vs.~Positive) & +0.184 (0.012) & +0.103 (0.076) &
+0.076 (0.002) & +0.116 (0.003) & +0.124 (0.003) & +0.065 (0.004) &
+0.224 (0.005) & +0.177 (0.003) \\
Positive + Nutrition Info (vs.~Positive) & +0.179 (0.012) & +0.203
(0.066) & +0.094 (0.002) & +0.105 (0.003) & +0.146 (0.003) & +0.069
(0.004) & +0.229 (0.004) & +0.168 (0.004) \\
Positive + High In (vs.~Positive) & +0.119 (0.012) & +0.132 (0.067) &
+0.011 (0.002) & -0.038 (0.003) & +0.000 (0.003) & -0.002 (0.004) &
+0.023 (0.005) & +0.042 (0.003) \\
Spectrum 1-to-5 star label (vs.~Positive) & +0.108 (0.012) & +0.122
(0.070) & +0.032 (0.003) & +0.057 (0.003) & +0.038 (0.003) & +0.057
(0.004) & +0.109 (0.005) & +0.139 (0.003) \\
High In warning label (vs.~Positive) & +0.054 (0.012) & +0.036 (0.069) &
+0.067 (0.002) & +0.014 (0.003) & -0.088 (0.003) & -0.022 (0.004) &
-0.047 (0.005) & -0.010 (0.003) \\
\end{longtable}
}

\emph{Note.} Ordinary least squares (OLS) point estimates and HC1
heteroskedasticity-robust standard errors (in parentheses) for the 5
front-of-package labeling arms relative to the \texttt{Positive}
reference arm under respondent-persona conditioning (\emph{N} = 200
shoppers, \emph{n} = 1,200 choices), alongside full-sample human
estimates (\emph{N} = 5,633 shoppers, 33,798 choices; Grummon et al.
(2025)). Kev rows use prompts built from the audited survey items with
13-attribute respondent personas.

\nexttablesize{\scriptsize}
{\def\LTcaptype{apptbl}
\begin{longtable}[]{@{}
  >{\raggedright\arraybackslash}p{(\linewidth - 14\tabcolsep) * \real{0.1250}}
  >{\raggedright\arraybackslash}p{(\linewidth - 14\tabcolsep) * \real{0.1250}}
  >{\raggedleft\arraybackslash}p{(\linewidth - 14\tabcolsep) * \real{0.1250}}
  >{\raggedleft\arraybackslash}p{(\linewidth - 14\tabcolsep) * \real{0.1250}}
  >{\raggedleft\arraybackslash}p{(\linewidth - 14\tabcolsep) * \real{0.1250}}
  >{\raggedleft\arraybackslash}p{(\linewidth - 14\tabcolsep) * \real{0.1250}}
  >{\raggedleft\arraybackslash}p{(\linewidth - 14\tabcolsep) * \real{0.1250}}
  >{\raggedleft\arraybackslash}p{(\linewidth - 14\tabcolsep) * \real{0.1250}}@{}}
\caption{{\label{apptbl-grummon2025-personas}}Single-attribute choice
cross-entropy across seen and unseen shopper covariates}\tabularnewline
\toprule\noalign{}
\begin{minipage}[b]{\linewidth}\raggedright
Attribute group
\end{minipage} & \begin{minipage}[b]{\linewidth}\raggedright
Shopper attribute
\end{minipage} & \begin{minipage}[b]{\linewidth}\raggedleft
Kev 0.8B (base)
\end{minipage} & \begin{minipage}[b]{\linewidth}\raggedleft
Kev 0.8B (FT)
\end{minipage} & \begin{minipage}[b]{\linewidth}\raggedleft
Kev 4B (base)
\end{minipage} & \begin{minipage}[b]{\linewidth}\raggedleft
Kev 4B (FT)
\end{minipage} & \begin{minipage}[b]{\linewidth}\raggedleft
Kev 9B (base)
\end{minipage} & \begin{minipage}[b]{\linewidth}\raggedleft
Kev 9B (FT)
\end{minipage} \\
\midrule\noalign{}
\endfirsthead
\toprule\noalign{}
\begin{minipage}[b]{\linewidth}\raggedright
Attribute group
\end{minipage} & \begin{minipage}[b]{\linewidth}\raggedright
Shopper attribute
\end{minipage} & \begin{minipage}[b]{\linewidth}\raggedleft
Kev 0.8B (base)
\end{minipage} & \begin{minipage}[b]{\linewidth}\raggedleft
Kev 0.8B (FT)
\end{minipage} & \begin{minipage}[b]{\linewidth}\raggedleft
Kev 4B (base)
\end{minipage} & \begin{minipage}[b]{\linewidth}\raggedleft
Kev 4B (FT)
\end{minipage} & \begin{minipage}[b]{\linewidth}\raggedleft
Kev 9B (base)
\end{minipage} & \begin{minipage}[b]{\linewidth}\raggedleft
Kev 9B (FT)
\end{minipage} \\
\midrule\noalign{}
\endhead
\bottomrule\noalign{}
\endlastfoot
Seen in training & Sex & 1.5253 & 1.5472 & 1.3948 & 1.4341 & 1.4902 &
1.3468 \\
Seen in training & Age & 1.4806 & 1.5426 & 1.3974 & 1.4366 & 1.5011 &
1.3532 \\
Seen in training & Education & 1.5152 & 1.5475 & 1.3897 & 1.4267 &
1.5375 & 1.3478 \\
Seen in training & Income & 1.4764 & 1.5214 & 1.4416 & 1.4362 & 1.4908 &
1.3462 \\
Seen in training & Race & 1.4746 & 1.5389 & 1.4243 & 1.4258 & 1.5054 &
1.3437 \\
Seen in training & How many people, including yourself, live in your
household? & 1.5395 & 1.5477 & 1.5033 & 1.4343 & 1.3946 & 1.3486 \\
\textbf{Seen mean} & --- & \textbf{1.5019} (6) & \textbf{1.5409} (6) &
\textbf{1.4252} (6) & \textbf{1.4323} (6) & \textbf{1.4866} (6) &
\textbf{1.3477} (6) \\
Unseen in training & Hispanic or Latino & 1.5420 & 1.5527 & 1.3834 &
1.4314 & 1.5001 & 1.3497 \\
Unseen in training & How many children (ages 0-18) currently live in
your household? & 1.5658 & 1.5499 & 1.5430 & 1.4370 & 1.3935 & 1.3420 \\
Unseen in training & SNAP / Food Stamps recipient & 1.5709 & 1.5419 &
1.4581 & 1.4295 & 1.5669 & 1.3507 \\
Unseen in training & Do you or anyone else in your household currently
receive WIC benefits? & 1.4839 & 1.5286 & 1.5181 & 1.4360 & 1.4193 &
1.3404 \\
Unseen in training & Nutrition Facts label usage & 1.4756 & 1.5366 &
1.4982 & 1.4372 & 1.6661 & 1.3576 \\
Unseen in training & Share of groceries bought online & 1.4754 & 1.5311
& 1.4355 & 1.4494 & 1.6016 & 1.3518 \\
Unseen in training & Nutrition literacy score & 1.4679 & 1.5418 & 1.4095
& 1.4387 & 1.5824 & 1.3391 \\
\textbf{Unseen mean} & --- & \textbf{1.5117} (7) & \textbf{1.5404} (7) &
\textbf{1.4637} (7) & \textbf{1.4370} (7) & \textbf{1.5328} (7) &
\textbf{1.3473} (7) \\
\end{longtable}
}

\emph{Note.} Three-option choice cross-entropy (in bits; uniform
baseline \(\log_2 3 = 1.5850\) bits) on the evaluation subsample
(\emph{N} = 200 shoppers, \emph{n} = 1,200 choices) of Grummon et al.
(2025) when conditioning each model on a single shopper attribute at a
time, grouped into demographic covariates present in the survey training
corpus and food-shopping and nutrition covariates unseen during
training. Kev runs cover 6 seen and 7 unseen attributes. Group means
average over the attributes in each group (count in parentheses).

\begin{appfig}

\centering{

\includegraphics[width=0.9\linewidth,height=\textheight,keepaspectratio]{figures/grummon2025_ols_random.png}

}

\caption{\label{appfig-grummon2025-random}OLS treatment-effect
estimation error on the nutrition labeling trial under random-persona
prompts}

\floatnote{

\emph{Note.} Panel A reports Mean Absolute Percentage Error (MAPE) and
Panel B reports Mean Squared Error (MSE) for participant-level HC1 OLS
treatment effects across the \(d = 5\) front-of-package labeling arms
(relative to \texttt{Positive}) under unconditioned random-persona
prompts (\emph{N} = 200 shoppers, \emph{n} = 1,200 choices) relative to
full-sample human \(\boldsymbol{\beta}^*\) (\emph{N} = 5,633; Grummon et
al. (2025)).

}

\end{appfig}%

\begin{appfig}

\centering{

\includegraphics[width=0.9\linewidth,height=\textheight,keepaspectratio]{figures/grummon2025_ols_persona.png}

}

\caption{\label{appfig-grummon2025-persona}OLS treatment-effect
estimation error on the nutrition labeling trial under
respondent-persona conditioning}

\floatnote{

\emph{Note.} Panel A reports Mean Absolute Percentage Error (MAPE) and
Panel B reports Mean Squared Error (MSE) for participant-level HC1 OLS
treatment effects across the \(d = 5\) front-of-package labeling arms
(relative to \texttt{Positive}) under 13-attribute respondent personas
(\emph{N} = 200 shoppers, \emph{n} = 1,200 choices) relative to
full-sample human \(\boldsymbol{\beta}^*\) (\emph{N} = 5,633; Grummon et
al. (2025)).

}

\end{appfig}%




\end{document}
